\chapter{Konstruktion und Beweis}
\section{Der Rekursionskern}

Diese Ergänzung enthält die Schnittkonstruktion und den vollständigen
Beweis des \FormulaRefAuto{DedekindRecursionTheorem}[thm] aus Band~10.
Ihre eigenen Aussagen tragen den Nummernpräfix \(10\mathrm E\).
Die \DedekindReadingLink{} entwickelt denselben Beweis zusammenhängend.


\paragraph{Beweisskizze}

\begingroup
\footnotesize
\[
\begin{array}{@{}c@{}}
\boxed{
\begin{array}{c}
m_0\in M,\qquad f\colon M\to M
\\[-.1ex]
\text{gesucht: eine eindeutig bestimmte Abbildung }
F\colon\mathbb{N}\to M
\\[-.1ex]
F(0)=m_0,\qquad F(\Succ(n))=f(F(n))
\end{array}}
\\[1.1ex]
\Downarrow
\\[-.2ex]
\boxed{
\begin{array}{c}
\AdmRec{m_0}{f}(H)
\\[-.1ex]
H\subseteq\mathbb{N}\times M,\qquad
(0,m_0)\in H
\\[-.1ex]
(n,a)\in H\Longrightarrow(\Succ(n),f(a))\in H
\end{array}}
\\[1.1ex]
\Downarrow
\\[-.2ex]
\boxed{
\begin{array}{c}
\RecCore{m_0}{f}
=
\displaystyle\bigcap\{H\subseteq\mathbb{N}\times M
\mid \AdmRec{m_0}{f}(H)\}
\\[-.1ex]
\text{kleinste rekursionsadmissible Teilmenge von }
\mathbb{N}\times M
\end{array}}
\\[1.1ex]
\Downarrow
\\[-.2ex]
\boxed{
\begin{array}{c}
\text{Nullstufe: }(0,a)\in\RecCore{m_0}{f}
\Longrightarrow a=m_0
\\[-.1ex]
\text{Nachfolgerstufe: }(n,a)\in\RecCore{m_0}{f}
\Longrightarrow
(\Succ(n),f(a))\in\RecCore{m_0}{f}
\end{array}}
\\[1.1ex]
\Downarrow
\\[-.2ex]
\boxed{
\begin{array}{c}
\forall n\in\mathbb{N}\,\exists!a\in M\,
\bigl((n,a)\in\RecCore{m_0}{f}\bigr)
\\[-.1ex]
\text{Existenz durch Induktion, Eindeutigkeit durch Induktion}
\end{array}}
\\[1.1ex]
\Downarrow
\\[-.2ex]
\boxed{
\begin{array}{c}
\RecFun{m_0}{f}\colon\mathbb{N}\to M,
\qquad
\RecFun{m_0}{f}(n)
=
\iota a\,\bigl((n,a)\in\RecCore{m_0}{f}\bigr)
\\[-.1ex]
\RecFun{m_0}{f}(0)=m_0,\qquad
\RecFun{m_0}{f}(\Succ(n))=f(\RecFun{m_0}{f}(n))
\end{array}}
\\[1.1ex]
\Downarrow
\\[-.2ex]
\boxed{
\begin{array}{c}
G\colon\mathbb{N}\to M,\quad
G(0)=m_0,\quad
G(\Succ(n))=f(G(n))
\\[-.1ex]
\Longrightarrow\quad
G=\RecFun{m_0}{f}
\end{array}}
\end{array}
\]
\endgroup

\par\Needspace{5\baselineskip}
Der Beweis ersetzt die direkte Definition einer Folge durch eine
Schnittkonstruktion. Zuerst werden diejenigen Teilmengen von
\(\mathbb{N}\times M\) betrachtet, die den Anfangswert enthalten und unter dem
Rekursionsschritt abgeschlossen sind. Ihr Durchschnitt ist der Rekursionskern.
Die Fixierungslemmata zeigen, dass der Kern in jeder Stufe höchstens einen
Wert enthält; die Induktion liefert zugleich die Existenz eines solchen Werts.
Damit ist der Kern der Graph einer Funktion. Die Rekursionsgleichungen folgen
dann aus der Konstruktion des Kerns, und die Eindeutigkeit einer beliebigen
anderen rekursiv definierten Abbildung wird erneut per Induktion bewiesen.

\paragraph{Aufbau des Abschnitts}

Der Abschnitt trennt die Hauptlinie des Beweises von den technischen
Fixierungsargumenten. Zuerst werden rekursionsadmissible Teilmengen und der
Rekursionskern konstruiert. Danach zeigen die Fixierungslemmata, dass jede
Stufe des Kerns höchstens einen Wert enthält; die Induktion liefert die
Existenz eines Werts in jeder Stufe. Aus eindeutiger Existenz wird der Kern als
Funktion gelesen. Die Rekursionsgleichungen und die Eindeutigkeit beliebiger
rekursiver Abbildungen ergeben schließlich den Dedekindschen Rekursionssatz.
Die in Band~10 definierte Rekursionsabbildung bezeichnet das durch diesen
Satz eindeutig bestimmte Objekt.

\subsubsection{Hauptkonstruktion: Rekursionsadmissible Mengen und Rekursionskern}

\paragraph{Axiome der Rekursionsadmissibilität}

\par\Needspace{14\baselineskip}
\FormulaAxiomDelta[Teilmenge des Produktraums]{%
  H\subseteq \mathbb{N}\times M
}{%
  \DeltaRow{Mengen}{H\dsep M}%
}

\par\Needspace{14\baselineskip}
\FormulaAxiomDelta[Startpaar]{%
  m_0\in M \vdash (0,m_0)\in H
}{%
  \DeltaRow{Mengen}{H\dsep M\dsep m_0}%
}

\par\Needspace{14\baselineskip}
\FormulaAxiomDelta[Rekursionsschritt]{%
  f\colon M\to M\dsep n\in\mathbb{N}\dsep a\in M\dsep (n,a)\in H
  \vdash (\Succ(n),f(a))\in H
}{%
  \DeltaRow{Mengen}{H\dsep M\dsep n\dsep a}%
  \DeltaRow{Funktionen}{f}%
}

\paragraph{Definition der Rekursionsadmissibilität}

\par\Needspace{14\baselineskip}
\FormulaDefDeltaR[Begriff der rekursionsadmissiblen Teilmenge]{%
  \AdmRec{m_0}{f}(H)
}{%
  \AdmRec{m_0}{f}(H)
}{%
  \DeltaRow{Mengen}{H\dsep M\dsep n\dsep a\dsep m_0}%
  \DeltaRow{Funktionen}{f}%
  \DeltaRow{\textbf{Axiome}}{}%

  \DeltaRow{\makecell[l]{Teilmenge des\\Produktraums}}
           {H\subseteq \mathbb{N}\times M}
           [\FormulaRefAuto{H\subseteq \mathbb{N}\times M}[ax]]

  \DeltaRow{Startpaar}
           {m_0\in M \vdash (0,m_0)\in H}
           [\FormulaRefAuto{m_0\in M \vdash (0,m_0)\in H}[ax]]

  \DeltaRow{\makecell[l]{Rekursions-\\schritt}}
           {%
             \begin{aligned}[t]
               &f\colon M\to M\dsep n\in\mathbb{N}\dsep{}\\
               &a\in M\dsep (n,a)\in H{}\\
               &\vdash (\Succ(n),f(a))\in H
             \end{aligned}
           }
           [\FormulaRefAuto{f\colon M\to M\dsep n\in\mathbb{N}\dsep a\in M\dsep (n,a)\in H \vdash (\Succ(n),f(a))\in H}[ax]]

  \DeltaRow{\textbf{Neue Symbole}}{}%
  \DeltaPrem{\makecell[l]{Rekursionsadmissible\\Teilmenge}}{ }%
}

\paragraph{Nichtleere Admissibilität und Hilfskonstruktionen}

Der volle Produktraum zeigt zunächst, dass es überhaupt rekursionsadmissible
Teilmengen gibt. Die beiden Fixierungsmengen sind technische Werkzeuge: Sie
werden später benutzt, um falsche Werte in der Nullstufe und in einer
Nachfolgerstufe aus dem Rekursionskern auszuschließen.

\subparagraph{Der volle Produktraum}

\par\Needspace{14\baselineskip}
\FormulaThmDelta[Der volle Produktraum ist rekursionsadmissibel]{%
  m_0\in M \dsep f\colon M\to M \vdash \AdmRec{m_0}{f}(\mathbb{N}\times M)
}{%
  \DeltaDecl{Mengen}{M\dsep n\dsep a\dsep m_0}%
  \DeltaDecl{Funktionen}{f}%
}
\begin{tabproofsplitwide}
  \proofpartwideind[Teilmenge des Produktraums]{%
    \mathbb{N}\times M \subseteq \mathbb{N}\times M}
    \proofstepwidestar{\mathbb{N}\times M \subseteq \mathbb{N}\times M}{%
      \FormulaRefAuto{A \subseteq A}}
  \closeproofpartwideind

  \proofpartwideindR[Startpaar]{%
    m_0\in M \vdash (0,m_0)\in \mathbb{N}\times M
  }{%
    m_0\in M \vdash (0,m_0)\in \mathbb{N}\times M
  }
    \proofstepwidestar[1]{m_0\in M}{\rA}
    \proofstepwidestar{0\in\mathbb{N}}{%
      \FormulaRefAuto{PeanoZero}}
    \proofstepwidestar[1]{(0,m_0)\in\mathbb{N}\times M}{%
      \FormulaRefAuto{a\in A \dsep b\in B \vdash (a,b)\in A\times B}{2,1}}
  \closeproofpartwideindR

  \proofpartwideindR[Rekursionsschritt]{%
    f\colon M\to M\dsep n\in\mathbb{N}\dsep a\in M\dsep
    (n,a)\in \mathbb{N}\times M \vdash
    (\Succ(n),f(a))\in \mathbb{N}\times M
  }{%
    f\colon M\to M\dsep n\in\mathbb{N}\dsep a\in M\dsep
    (n,a)\in \mathbb{N}\times M \vdash
    (\Succ(n),f(a))\in \mathbb{N}\times M
  }
    \proofstepwidestar[1]{f\colon M\to M}{\rA}
    \proofstepwidestar[2]{n\in\mathbb{N}}{\rA}
    \proofstepwidestar[3]{a\in M}{\rA}
    \proofstepwidestar[4]{(n,a)\in \mathbb{N}\times M}{\rA}

    \proofstepwidestar[2]{\Succ(n)\in\mathbb{N}}{%
      \FormulaRefAuto{PeanoSuccClosure}{2}}

    \proofstepwidestar[1,3]{f(a)\in M}{%
      \FormulaRefAuto{F\colon A\to B\dsep x\in A\vdash F(x)\in B}{1,3}}

    \proofstepwidestar[1,2,3,4]{(\Succ(n),f(a))\in\mathbb{N}\times M}{%
      \FormulaRefAuto{a\in A \dsep b\in B \vdash (a,b)\in A\times B}{5,6}}
  \closeproofpartwideindR

  \proofpartwide{Schluss}
    \proofstepwidestar[1]{m_0\in M}{\rA}
    \proofstepwidestar[2]{f\colon M\to M}{\rA}

    \proofstepwidestar[]{\mathbb{N}\times M \subseteq \mathbb{N}\times M}{%
      \FormulaRefAuto{\mathbb{N}\times M \subseteq \mathbb{N}\times M}}
    \proofstepwidestar[1]{(0,m_0)\in \mathbb{N}\times M}{%
      \FormulaRefAuto{m_0\in M \vdash (0,m_0)\in \mathbb{N}\times M}{1}}
    \proofstepwidestar[2]{\forall n\in\mathbb N\,\forall a\in M\,\bigl((n,a)\in \mathbb{N}\times M\rightarrow (\Succ(n),f(a))\in \mathbb{N}\times M\bigr)}{%
      \FormulaRefAuto{f\colon M\to M\dsep n\in\mathbb{N}\dsep a\in M\dsep (n,a)\in \mathbb{N}\times M \vdash (\Succ(n),f(a))\in \mathbb{N}\times M}{2}}
    \proofstepwidestar[1,2]{\AdmRec{m_0}{f}(\mathbb{N}\times M)}{\FormulaRefAuto{\AdmRec{m_0}{f}(H)}{3,4,5}}
  \closeproofpartwide
\end{tabproofsplitwide}

\subparagraph{Technisches Lemma: Fixierungsmenge der Nullstufe}

\par\Needspace{14\baselineskip}
\FormulaDefDeltaR[Fixierungsmenge der Nullstufe]{%
  \begin{aligned}[t]
    &m_0\in M \vdash{}\\
    &\ZeroFix{m_0}\coloneq
    \{(n,a)\in \mathbb{N}\times M \mid n=0 \rightarrow a=m_0\}
  \end{aligned}
}{%
  m_0\in M \vdash
  \ZeroFix{m_0}\coloneq
  \{(n,a)\in \mathbb{N}\times M \mid n=0 \rightarrow a=m_0\} }{%
  \DeltaDecl{Mengen}{M\dsep n\dsep a\dsep m_0}%
}

\par\Needspace{14\baselineskip}
\FormulaThmDelta[Fixierungsmenge der Nullstufe ist rekursionsadmissibel]{%
  m_0\in M\dsep f\colon M\to M \vdash \AdmRec{m_0}{f}(\ZeroFix{m_0}) }{%
  \DeltaDecl{Mengen}{M\dsep n\dsep a\dsep m_0}%
  \DeltaDecl{Funktionen}{f}%
}
\begin{tabproofsplitwide}
  \proofpartwideindR[Teilmenge des Produktraums]{%
    m_0\in M \vdash \ZeroFix{m_0}\subseteq \mathbb{N}\times M
  }{%
    m_0\in M \vdash \ZeroFix{m_0}\subseteq \mathbb{N}\times M
  }
    \proofstepwidestar[1]{m_0\in M}{\rA}
    \proofstepwidestar[]{%
      \{(n,a)\in \mathbb{N}\times M \mid n=0 \rightarrow a=m_0\}
      \subseteq \mathbb{N}\times M
}{%
      \FormulaRefAuto{\{ x \in A \mid P(x) \} \subseteq A}}
    \proofstepwidestar[1]{\ZeroFix{m_0}\subseteq \mathbb{N}\times M}{%
      \FormulaRefAuto{m_0\in M \vdash \ZeroFix{m_0}\coloneq \{(n,a)\in \mathbb{N}\times M \mid n=0 \rightarrow a=m_0\}}{1,2}}
  \closeproofpartwideindR

  \proofpartwideindR[Startpaar]{%
    m_0\in M \vdash (0,m_0)\in \ZeroFix{m_0}
  }{%
    m_0\in M \vdash (0,m_0)\in \ZeroFix{m_0}
  }
    \proofstepwidestar[1]{m_0\in M}{\rA}
    \proofstepwidestar{0\in\mathbb{N}}{%
      \FormulaRefAuto{PeanoZero}}
    \proofstepwidestar[1]{(0,m_0)\in \mathbb{N}\times M}{%
      \FormulaRefAuto{a\in A \dsep b\in B \vdash (a,b)\in A\times B}{2,1}}
    \proofstepwidestar[]{m_0=m_0}{\rII}
    \proofstepwidestar[]{%
      0=0 \rightarrow m_0=m_0
    }{%
      \FormulaRefAuto{Q\vdash P \rightarrow Q}{4}}
    \proofstepwidestar[1]{%
      (0,m_0)\in
      \{(n,a)\in \mathbb{N}\times M \mid n=0 \rightarrow a=m_0\}
    }{%
      \FormulaRefAuto{x \in A\dsep P(x)\vdash x \in \{x \in A \mid P(x)\}}{3,5}}
    \proofstepwidestar[1]{(0,m_0)\in \ZeroFix{m_0}}{%
      \FormulaRefAuto{m_0\in M \vdash \ZeroFix{m_0}\coloneq \{(n,a)\in \mathbb{N}\times M \mid n=0 \rightarrow a=m_0\}}{1,6}}
  \closeproofpartwideindR

  \proofpartwideindR[Rekursionsschritt]{%
    m_0\in M\dsep f\colon M\to M\dsep n\in\mathbb{N}\dsep a\in M\dsep
    (n,a)\in \ZeroFix{m_0} \vdash
    (\Succ(n),f(a))\in \ZeroFix{m_0}
  }{%
    m_0\in M\dsep f\colon M\to M\dsep n\in\mathbb{N}\dsep a\in M\dsep
    (n,a)\in \ZeroFix{m_0} \vdash
    (\Succ(n),f(a))\in \ZeroFix{m_0}
  }
    \proofstepwidestar[1]{m_0\in M}{\rA}
    \proofstepwidestar[2]{f\colon M\to M}{\rA}
    \proofstepwidestar[3]{n\in\mathbb{N}}{\rA}
    \proofstepwidestar[4]{a\in M}{\rA}
    \proofstepwidestar[5]{(n,a)\in \ZeroFix{m_0}}{\rA}

    \proofstepwidestar[3]{\Succ(n)\in\mathbb{N}}{%
      \FormulaRefAuto{PeanoSuccClosure}{3}}

    \proofstepwidestar[2,4]{f(a)\in M}{%
      \FormulaRefAuto{F\colon A\to B\dsep x\in A\vdash F(x)\in B}{2,4}}

    \proofstepwidestar[2,3,4]{(\Succ(n),f(a))\in \mathbb{N}\times M}{%
      \FormulaRefAuto{a\in A \dsep b\in B \vdash (a,b)\in A\times B}{6,7}}

    \proofstepwidestar[3]{\Succ(n)\neq0}{%
      \FormulaRefAuto{PeanoSuccNonzero}{3}}

    \proofstepwidestar[3]{\Succ(n)=0 \rightarrow f(a)=m_0}{%
      \FormulaRefAuto{\neg P \vdash P \rightarrow Q}{9}}

    \proofstepwidestar[1,2,3,4,5]{%
      (\Succ(n),f(a))\in
      \{(n,a)\in \mathbb{N}\times M \mid n=0 \rightarrow a=m_0\}
    }{%
      \FormulaRefAuto{x \in A\dsep P(x)\vdash x \in \{x \in A \mid P(x)\}}{8,10}}

    \proofstepwidestar[1,2,3,4,5]{(\Succ(n),f(a))\in \ZeroFix{m_0}}{%
      \FormulaRefAuto{m_0\in M \vdash \ZeroFix{m_0}\coloneq \{(n,a)\in \mathbb{N}\times M \mid n=0 \rightarrow a=m_0\}}{1,11}}
  \closeproofpartwideindR

  \proofpartwide{Schluss}
    \proofstepwidestar[1]{m_0\in M}{\rA}
    \proofstepwidestar[2]{f\colon M\to M}{\rA}

    \proofstepwidestar[1]{(0,m_0)\in \ZeroFix{m_0}}{%
      \FormulaRefAuto{m_0\in M \vdash (0,m_0)\in \ZeroFix{m_0}}{1}}
    \proofstepwidestar[1]{\ZeroFix{m_0}\subseteq \mathbb{N}\times M}{%
      \FormulaRefAuto{m_0\in M \vdash \ZeroFix{m_0}\subseteq \mathbb{N}\times M}{1}}
    \proofstepwidestar[1,2]{\forall n\in\mathbb N\,\forall a\in M\,\bigl((n,a)\in \ZeroFix{m_0}\rightarrow (\Succ(n),f(a))\in \ZeroFix{m_0}\bigr)}{%
      \FormulaRefAuto{m_0\in M\dsep f\colon M\to M\dsep n\in\mathbb{N}\dsep a\in M\dsep (n,a)\in \ZeroFix{m_0} \vdash (\Succ(n),f(a))\in \ZeroFix{m_0}}{1,2}}
    \proofstepwidestar[1,2]{\AdmRec{m_0}{f}(\ZeroFix{m_0})}{\FormulaRefAuto{\AdmRec{m_0}{f}(H)}{4,3,5}}
  \closeproofpartwide
\end{tabproofsplitwide}

\subparagraph{Schnittkonstruktion des Rekursionskerns}

Der Rekursionskern ist die kleinste rekursionsadmissible Teilmenge von
\(\mathbb{N}\times M\): Ein Paar liegt genau dann im Kern, wenn es in jeder
rekursionsadmissiblen Teilmenge liegt.

\par\Needspace{14\baselineskip}
\FormulaDefDelta[Rekursionskern]{%
  m_0\in M \dsep f\colon M\to M \vdash
  \RecCore{m_0}{f}\coloneq
  \bigcap_{\AdmRec{m_0}{f}(H)} H
}{%
  \DeltaDecl{Mengen}{H\dsep M\dsep m_0}%
  \DeltaDecl{Funktionen}{f}%
}

\par\Needspace{14\baselineskip}
\FormulaThmDeltaR[Charakterisierung des Rekursionskerns]{%
  \begin{aligned}
    &m_0\in M\dsep f\colon M\to M \vdash{}\\
    &(n,a)\in \RecCore{m_0}{f}\leftrightarrow{}\\
    &\forall H\,\bigl(\AdmRec{m_0}{f}(H)\rightarrow (n,a)\in H\bigr)
  \end{aligned}
}{%
  m_0\in M\dsep f\colon M\to M \vdash
  (n,a)\in \RecCore{m_0}{f}\leftrightarrow
  \forall H\,\bigl(\AdmRec{m_0}{f}(H)\rightarrow (n,a)\in H\bigr)
}{%
  \DeltaDecl{Mengen}{H\dsep M\dsep n\dsep a\dsep m_0}%
  \DeltaDecl{Funktionen}{f}%
}
\begin{tabproofsplitwide}
  \proofpartwideindR[Hinrichtung]{%
    \begin{aligned}[t]
      &m_0\in M\dsep f\colon M\to M \vdash{}\\
      &(n,a)\in \RecCore{m_0}{f}\rightarrow{}\\
      &\forall H\,\bigl(\AdmRec{m_0}{f}(H)\rightarrow (n,a)\in H\bigr)
    \end{aligned}
  }{%
    m_0\in M\dsep f\colon M\to M\dsep (n,a)\in \RecCore{m_0}{f}
    \vdash \forall H\,\bigl(\AdmRec{m_0}{f}(H)\rightarrow (n,a)\in H\bigr)
  }
    \proofstepwidestar[1]{m_0\in M}{\rA}
    \proofstepwidestar[2]{f\colon M\to M}{\rA}
    \proofstepwidestar[3]{(n,a)\in \RecCore{m_0}{f}}{\rA}

    \proofstepwidestar[1,2]{\AdmRec{m_0}{f}(\mathbb{N}\times M)}{%
      \FormulaRefAuto{m_0\in M\dsep f\colon M\to M\vdash \AdmRec{m_0}{f}(\mathbb{N}\times M)}{1,2}}
    \proofstepwidestar[1,2]{\exists H\,\AdmRec{m_0}{f}(H)}{\rEI{4}}

    \proofstepwidestar[1,2]{\RecCore{m_0}{f}=\bigcap_{\AdmRec{m_0}{f}(H)}H}{%
      \FormulaRefAuto{m_0\in M\dsep f\colon M\to M\vdash
        \RecCore{m_0}{f}\coloneq \bigcap_{\AdmRec{m_0}{f}(H)} H}{1,2}}
    \proofstepwidestar[1,2,3]{(n,a)\in \bigcap_{\AdmRec{m_0}{f}(H)} H}{%
      \rIE{6,3}}

    \proofstepwidestar[1,2,3]{%
      \forall H\,\bigl(\AdmRec{m_0}{f}(H)\rightarrow (n,a)\in H\bigr)
    }{%
      \FormulaRefAuto{%
        \exists A(P(A)) \vdash
        x \in \bigcap_{P(B)} B \leftrightarrow
        \forall C\, (P(C) \rightarrow x \in C)
      }{5,7}}

    \proofstepwidestar[1,2]{%
    \begin{aligned}[t]
      &(n,a)\in \RecCore{m_0}{f}\rightarrow{}\\
      &\forall H\,\bigl(\AdmRec{m_0}{f}(H)\rightarrow (n,a)\in H\bigr)
    \end{aligned}
    }{\rRI{3,8}}
  \closeproofpartwideindR

  \proofpartwideindR[Rückrichtung]{%
    \begin{aligned}[t]
      &m_0\in M\dsep f\colon M\to M \vdash{}\\
      &\forall H\,\bigl(\AdmRec{m_0}{f}(H)\rightarrow (n,a)\in H\bigr)\rightarrow{}\\
      &(n,a)\in \RecCore{m_0}{f}
    \end{aligned}
  }{%
    m_0\in M\dsep f\colon M\to M\dsep
    \forall H\,\bigl(\AdmRec{m_0}{f}(H)\rightarrow (n,a)\in H\bigr)
    \vdash (n,a)\in \RecCore{m_0}{f}
  }
    \proofstepwidestar[1]{m_0\in M}{\rA}
    \proofstepwidestar[2]{f\colon M\to M}{\rA}
    \proofstepwidestar[3]{%
      \forall H\,\bigl(\AdmRec{m_0}{f}(H)\rightarrow (n,a)\in H\bigr)
    }{\rA}

    \proofstepwidestar[1,2]{\AdmRec{m_0}{f}(\mathbb{N}\times M)}{%
      \FormulaRefAuto{m_0\in M\dsep f\colon M\to M\vdash \AdmRec{m_0}{f}(\mathbb{N}\times M)}{1,2}}
    \proofstepwidestar[1,2]{\exists H\,\AdmRec{m_0}{f}(H)}{\rEI{4}}

    \proofstepwidestar[1,2,3]{(n,a)\in \bigcap_{\AdmRec{m_0}{f}(H)} H}{%
      \FormulaRefAuto{%
        \exists A(P(A)) \vdash
        x \in \bigcap_{P(B)} B \leftrightarrow
        \forall C\, (P(C) \rightarrow x \in C)
      }{5,3}}

    \proofstepwidestar[1,2]{\RecCore{m_0}{f}=\bigcap_{\AdmRec{m_0}{f}(H)}H}{%
      \FormulaRefAuto{m_0\in M\dsep f\colon M\to M\vdash
        \RecCore{m_0}{f}\coloneq \bigcap_{\AdmRec{m_0}{f}(H)} H}{1,2}}
    \proofstepwidestar[1,2,3]{(n,a)\in \RecCore{m_0}{f}}{%
      \rIE{7,6}}

    \proofstepwidestar[1,2]{%
    \begin{aligned}[t]
      &\forall H\,\bigl(\AdmRec{m_0}{f}(H)\rightarrow (n,a)\in H\bigr)\rightarrow{}\\
      &(n,a)\in \RecCore{m_0}{f}
    \end{aligned}
    }{\rRI{3,8}}
  \closeproofpartwideindR

\proofpartwide{Schluss}
  \proofstepwidestar[1]{m_0\in M}{\rA}
  \proofstepwidestar[2]{f\colon M\to M}{\rA}

  \proofstepwidestar[1,2]{%
    \begin{aligned}[t]
      &(n,a)\in \RecCore{m_0}{f}\rightarrow{}\\
      &\forall H\,\bigl(\AdmRec{m_0}{f}(H)\rightarrow (n,a)\in H\bigr)
    \end{aligned}
  }{%
    \FormulaRefAuto{%
      m_0\in M\dsep f\colon M\to M\dsep (n,a)\in \RecCore{m_0}{f}
      \vdash \forall H\,\bigl(\AdmRec{m_0}{f}(H)\rightarrow (n,a)\in H\bigr)
    }{1,2}}

  \proofstepwidestar[1,2]{%
    \begin{aligned}[t]
      &\forall H\,\bigl(\AdmRec{m_0}{f}(H)\rightarrow (n,a)\in H\bigr)\rightarrow{}\\
      &(n,a)\in \RecCore{m_0}{f}
    \end{aligned}
  }{%
    \FormulaRefAuto{%
      m_0\in M\dsep f\colon M\to M\dsep
      \forall H\,\bigl(\AdmRec{m_0}{f}(H)\rightarrow (n,a)\in H\bigr)
      \vdash (n,a)\in \RecCore{m_0}{f}
    }{1,2}}

  \proofstepwidestar[1,2]{%
    \begin{aligned}[t]
      &(n,a)\in \RecCore{m_0}{f}\leftrightarrow{}\\
      &\forall H\,\bigl(\AdmRec{m_0}{f}(H)\rightarrow (n,a)\in H\bigr)
    \end{aligned}
  }{\rLRI{3,4}}
\closeproofpartwide
\end{tabproofsplitwide}

\par\Needspace{14\baselineskip}
\FormulaThmDelta[Minimalität des Rekursionskerns]{%
  m_0\in M\dsep f\colon M\to M\dsep \AdmRec{m_0}{f}(H)
  \vdash \RecCore{m_0}{f}\subseteq H
}{%
  \DeltaDecl{Mengen}{H\dsep M\dsep m_0}%
  \DeltaDecl{Funktionen}{f}%
}
\begin{tabproof}
  \proofstep{1}{m_0\in M}{\rA}
  \proofstep{2}{f\colon M\to M}{\rA}
  \proofstep{3}{\AdmRec{m_0}{f}(H)}{\rA}
  \proofstep{3}{\bigcap_{\AdmRec{m_0}{f}(X)} X \subseteq H}{%
    \FormulaRefAuto{P(C)\vdash \bigcap_{P(A)} A \subseteq C}{3}}
  \proofstep{1,2}{\RecCore{m_0}{f}=\bigcap_{\AdmRec{m_0}{f}(H)}H}{%
      \FormulaRefAuto{%
      m_0\in M \dsep f\colon M\to M \vdash
      \RecCore{m_0}{f}\coloneq
      \bigcap_{\AdmRec{m_0}{f}(X)} X
    }{1,2}}
  \proofstep{1,2,3}{\RecCore{m_0}{f}\subseteq H}{%
    \rIE{5,4}}
\end{tabproof}

\par\Needspace{14\baselineskip}
\FormulaThmDelta[Der Rekursionskern liegt im Produktraum]{%
  m_0\in M\dsep f\colon M\to M \vdash \RecCore{m_0}{f}\subseteq \mathbb{N}\times M
}{%
  \DeltaDecl{Mengen}{M\dsep m_0}%
  \DeltaDecl{Funktionen}{f}%
}
\begin{tabproof}
  \proofstep{1}{m_0\in M}{\rA}
  \proofstep{2}{f\colon M\to M}{\rA}
  \proofstep{1,2}{\AdmRec{m_0}{f}(\mathbb{N}\times M)}{%
    \FormulaRefAuto{m_0\in M \dsep f\colon M\to M \vdash \AdmRec{m_0}{f}(\mathbb{N}\times M)}{1,2}}
  \proofstep{1,2}{\RecCore{m_0}{f}\subseteq \mathbb{N}\times M}{%
    \FormulaRefAuto{m_0\in M\dsep f\colon M\to M\dsep \AdmRec{m_0}{f}(H)
    \vdash \RecCore{m_0}{f}\subseteq H}{1,2,3}}
\end{tabproof}

\par\Needspace{14\baselineskip}
\FormulaThmDelta[Das Startpaar liegt im Rekursionskern]{%
  m_0\in M\dsep f\colon M\to M \vdash (0,m_0)\in \RecCore{m_0}{f}
}{%
  \DeltaDecl{Mengen}{H\dsep M\dsep m_0}%
  \DeltaDecl{Funktionen}{f}%
}
\begin{tabproof}
  \proofstep{1}{m_0\in M}{\rA}
  \proofstep{2}{f\colon M\to M}{\rA}

  \proofstep{3}{\AdmRec{m_0}{f}(H)}{\rA}
  \proofstep{1,3}{(0,m_0)\in H}{%
    \FormulaRefAuto{\AdmRec{m_0}{f}(H)}{3}}

  \proofstep{}{%
    \forall H\,\bigl(\AdmRec{m_0}{f}(H)\rightarrow (0,m_0)\in H\bigr)
  }{\rUI{\rRI{3,4}}}

  \proofstep{1,2}{(0,m_0)\in \RecCore{m_0}{f}}{%
    \FormulaRefAuto{%
      m_0\in M\dsep f\colon M\to M \vdash
      (n,a)\in \RecCore{m_0}{f}\leftrightarrow
      \forall H\,\bigl(\AdmRec{m_0}{f}(H)\rightarrow (n,a)\in H\bigr)
    }{1,2,5}}
\end{tabproof}

\par\Needspace{14\baselineskip}
\FormulaThmDeltaR[Abschluss des Rekursionskerns]{%
  \begin{aligned}
    &m_0\in M\dsep f\colon M\to M\dsep n\in\mathbb{N}\dsep{}\\
    &a\in M\dsep (n,a)\in \RecCore{m_0}{f}
      \vdash (\Succ(n),f(a))\in \RecCore{m_0}{f}
  \end{aligned}
}{%
  m_0\in M\dsep f\colon M\to M\dsep n\in\mathbb{N}\dsep
  a\in M\dsep (n,a)\in \RecCore{m_0}{f}
  \vdash (\Succ(n),f(a))\in \RecCore{m_0}{f}
}{%
  \DeltaDecl{Mengen}{H\dsep M\dsep n\dsep a\dsep m_0}%
  \DeltaDecl{Funktionen}{f}%
}
\begin{tabproof}
  \proofstep{1}{m_0\in M}{\rA}
  \proofstep{2}{f\colon M\to M}{\rA}
  \proofstep{3}{n\in\mathbb{N}}{\rA}
  \proofstep{4}{a\in M}{\rA}
  \proofstep{5}{(n,a)\in \RecCore{m_0}{f}}{\rA}

  \proofstep{1,2}{%
    \begin{aligned}[t]
      &(n,a)\in \RecCore{m_0}{f}\leftrightarrow{}\\
      &\forall H\,\bigl(\AdmRec{m_0}{f}(H)\rightarrow (n,a)\in H\bigr)
    \end{aligned}
  }{%
    \FormulaRefAuto{%
      m_0\in M\dsep f\colon M\to M \vdash
      (n,a)\in \RecCore{m_0}{f}\leftrightarrow
      \forall H\,\bigl(\AdmRec{m_0}{f}(H)\rightarrow (n,a)\in H\bigr)
    }{1,2}}

  \proofstep{1,2,5}{%
    \begin{aligned}[t]
      &\forall H\,\bigl(\AdmRec{m_0}{f}(H)\rightarrow{}\\
      &\qquad (n,a)\in H\bigr)
    \end{aligned}
  }{%
    \FormulaRefAuto{P \leftrightarrow Q, P \vdash Q}{6,5}}

  \proofstep{8}{\AdmRec{m_0}{f}(H)}{\rA}
  \proofstep{1,2,5,8}{(n,a)\in H}{\rRE{\rUE{7},8}}

  \proofstep{2,3,4,8,9}{(\Succ(n),f(a))\in H}{%
    \FormulaRefAuto{f\colon M\to M\dsep n\in\mathbb{N}\dsep a\in M\dsep (n,a)\in H \vdash (\Succ(n),f(a))\in H}[ax]{2,3,4,9}}

  \proofstep{1,2,3,4,5}{%
    \begin{aligned}[t]
      &\forall H\,\bigl(\AdmRec{m_0}{f}(H)\rightarrow{}\\
      &\qquad (\Succ(n),f(a))\in H\bigr)
    \end{aligned}
  }{\rUI{\rRI{8,10}}}

  \proofstep{1,2}{%
    \begin{aligned}[t]
      &(\Succ(n),f(a))\in \RecCore{m_0}{f}\leftrightarrow{}\\
      &\forall H\,\bigl(\AdmRec{m_0}{f}(H)\rightarrow (\Succ(n),f(a))\in H\bigr)
    \end{aligned}
  }{%
    \FormulaRefAuto{%
      m_0\in M\dsep f\colon M\to M \vdash
      (n,a)\in \RecCore{m_0}{f}\leftrightarrow
      \forall H\,\bigl(\AdmRec{m_0}{f}(H)\rightarrow (n,a)\in H\bigr)
    }{1,2}}

  \proofstep{1,2,3,4,5}{(\Succ(n),f(a))\in \RecCore{m_0}{f}}{%
    \FormulaRefAuto{P \leftrightarrow Q, Q \vdash P}{12,11}}
\end{tabproof}

\par\Needspace{14\baselineskip}
\FormulaThmDelta[Der Rekursionskern ist rekursionsadmissibel]{%
  m_0\in M\dsep f\colon M\to M \vdash \AdmRec{m_0}{f}(\RecCore{m_0}{f})
}{%
  \DeltaDecl{Mengen}{H\dsep M\dsep n\dsep a\dsep m_0}%
  \DeltaDecl{Funktionen}{f}%
}
\begin{tabproofwide}
  \proofstepwidestar[1]{m_0\in M}{\rA}
  \proofstepwidestar[2]{f\colon M\to M}{\rA}

  \proofstepwidestar[1,2]{(0,m_0)\in \RecCore{m_0}{f}}{%
      \FormulaRefAuto{m_0\in M\dsep f\colon M\to M \vdash (0,m_0)\in \RecCore{m_0}{f}}{1,2}}
  \proofstepwidestar[1,2]{\RecCore{m_0}{f}\subseteq \mathbb{N}\times M}{%
      \FormulaRefAuto{m_0\in M\dsep f\colon M\to M \vdash \RecCore{m_0}{f}\subseteq \mathbb{N}\times M}{1,2}}
  \proofstepwidestar[1,2]{\forall n\in\mathbb N\,\forall a\in M\,\bigl((n,a)\in \RecCore{m_0}{f}\rightarrow (\Succ(n),f(a))\in \RecCore{m_0}{f}\bigr)}{%
      \FormulaRefAuto{%
        m_0\in M\dsep f\colon M\to M\dsep n\in\mathbb{N}\dsep
        a\in M\dsep (n,a)\in \RecCore{m_0}{f}
        \vdash (\Succ(n),f(a))\in \RecCore{m_0}{f}
      }{1,2}}
  \proofstepwidestar[1,2]{\AdmRec{m_0}{f}(\RecCore{m_0}{f})}{\FormulaRefAuto{\AdmRec{m_0}{f}(H)}{4,3,5}}
\end{tabproofwide}

\par\Needspace{14\baselineskip}
\FormulaThmDeltaK[Rekursionskern ist kleinster rekursionsadmissibler Graphkandidat]{%
  \begin{aligned}[t]
    &m_0\in M\dsep f\colon M\to M \vdash{}\\
    &\RecCore{m_0}{f}\subseteq \mathbb{N}\times M \land{}\\
    &(0,m_0)\in \RecCore{m_0}{f}\land{}\\
    &\AdmRec{m_0}{f}(\RecCore{m_0}{f})\land{}\\
    &\forall H\,\bigl(\AdmRec{m_0}{f}(H)\rightarrow
      \RecCore{m_0}{f}\subseteq H\bigr)
  \end{aligned}
}{RecCoreLeastAdmissible}{%
  \DeltaDecl{Mengen}{H\dsep M\dsep m_0}%
  \DeltaDecl{Funktionen}{f}%
}
\begin{tabproofsplitwide}
\proofpartwide{Trägerinklusion}
\proofstepwidestar[1]{m_0\in M}{\rA}
  \proofstepwidestar[2]{f\colon M\to M}{\rA}

  \proofstepwidestar[1,2]{\RecCore{m_0}{f}\subseteq \mathbb{N}\times M}{%
    \FormulaRefAuto{m_0\in M\dsep f\colon M\to M \vdash \RecCore{m_0}{f}\subseteq \mathbb{N}\times M}{1,2}}

\closeproofpartwide

\proofpartwide{Startpaar}
  \setcounter{proofstepnr}{3}% Fortlaufende Schritte dieses Beweises.
\proofstepwidestar[1,2]{(0,m_0)\in \RecCore{m_0}{f}}{%
    \FormulaRefAuto{m_0\in M\dsep f\colon M\to M \vdash (0,m_0)\in \RecCore{m_0}{f}}{1,2}}

\closeproofpartwide

\proofpartwide{Rekursionsabschluss}
  \setcounter{proofstepnr}{4}% Fortlaufende Schritte dieses Beweises.
\proofstepwidestar[1,2]{\AdmRec{m_0}{f}(\RecCore{m_0}{f})}{%
    \FormulaRefAuto{m_0\in M\dsep f\colon M\to M \vdash \AdmRec{m_0}{f}(\RecCore{m_0}{f})}{1,2}}


\closeproofpartwide

\proofpartwide{Minimalität}
  \setcounter{proofstepnr}{5}% Fortlaufende Schritte dieses Beweises.
\proofstepwidestar[6]{\AdmRec{m_0}{f}(H)}{\rA}
  \proofstepwidestar[1,2,6]{\RecCore{m_0}{f}\subseteq H}{%
    \FormulaRefAuto{m_0\in M\dsep f\colon M\to M\dsep \AdmRec{m_0}{f}(H)
      \vdash \RecCore{m_0}{f}\subseteq H}{1,2,6}}
  \proofstepwidestar[1,2]{%
    \forall H\,\bigl(\AdmRec{m_0}{f}(H)\rightarrow
    \RecCore{m_0}{f}\subseteq H\bigr)
  }{\rUI{\rRI{6,7}}}


\closeproofpartwide

\proofpartwide{Zusammenführung}
  \setcounter{proofstepnr}{8}% Fortlaufende Schritte dieses Beweises.
\proofstepwidestar[1,2]{%
    \begin{aligned}[t]
      &\RecCore{m_0}{f}\subseteq \mathbb{N}\times M\land{}\\
      &(0,m_0)\in \RecCore{m_0}{f}
    \end{aligned}
  }{\rAI{3,4}}
  \proofstepwidestar[1,2]{%
    \begin{aligned}[t]
      &\RecCore{m_0}{f}\subseteq \mathbb{N}\times M\land{}\\
      &(0,m_0)\in \RecCore{m_0}{f}\land{}\\
      &\AdmRec{m_0}{f}(\RecCore{m_0}{f})
    \end{aligned}
  }{\rAI{9,5}}
  \proofstepwidestar[1,2]{%
    \begin{aligned}[t]
      &\RecCore{m_0}{f}\subseteq \mathbb{N}\times M\land{}\\
      &(0,m_0)\in \RecCore{m_0}{f}\land{}\\
      &\AdmRec{m_0}{f}(\RecCore{m_0}{f})\land{}\\
      &\forall H\,\bigl(\AdmRec{m_0}{f}(H)\rightarrow
        \RecCore{m_0}{f}\subseteq H\bigr)
    \end{aligned}
  }{\rAI{10,8}}

\closeproofpartwide
\end{tabproofsplitwide}

\subparagraph{Fixierung der Nullstufe des Kerns}

\par\Needspace{14\baselineskip}
\FormulaThmDelta[Die Nullstufe des Rekursionskerns ist festgelegt]{%
  m_0\in M\dsep f\colon M\to M\dsep
  (0,a)\in \RecCore{m_0}{f}\vdash a=m_0
}{%
  \DeltaDecl{Mengen}{M\dsep a\dsep m_0}%
  \DeltaDecl{Funktionen}{f}%
}
\begin{tabproofwide}
  \proofstepwidestar[1]{m_0\in M}{\rA}
  \proofstepwidestar[2]{f\colon M\to M}{\rA}
  \proofstepwidestar[3]{(0,a)\in \RecCore{m_0}{f}}{\rA}

  \proofstepwidestar[1,2]{\AdmRec{m_0}{f}(\ZeroFix{m_0})}{%
    \FormulaRefAuto{m_0\in M\dsep f\colon M\to M \vdash \AdmRec{m_0}{f}(\ZeroFix{m_0})}{1,2}}

  \proofstepwidestar[1,2]{\RecCore{m_0}{f}\subseteq \ZeroFix{m_0}}{%
    \FormulaRefAuto{m_0\in M\dsep f\colon M\to M\dsep \AdmRec{m_0}{f}(H)\vdash \RecCore{m_0}{f}\subseteq H}{1,2,4}}

  \proofstepwidestar[1,2,3]{(0,a)\in \ZeroFix{m_0}}{%
    \FormulaRefAuto{A\subseteq B\dsep x\in A\vdash x\in B}{5,3}}

  \proofstepwidestar[1]{%
    \begin{aligned}[t]
      &\ZeroFix{m_0}\coloneq{}\\
      &\{(n,a)\in \mathbb{N}\times M \mid{}\\
      &\qquad n=0 \rightarrow a=m_0\}
    \end{aligned}
  }{%
    \FormulaRefAuto{m_0\in M \vdash \ZeroFix{m_0}\coloneq \{(n,a)\in \mathbb{N}\times M \mid n=0 \rightarrow a=m_0\}}{1}}

  \proofstepwidestar[1,3]{%
    (0,a)\in
    \{(n,a)\in \mathbb{N}\times M \mid n=0 \rightarrow a=m_0\}
  }{%
    \rIE{7,6}}

  \proofstepwidestar[1,3]{0=0 \rightarrow a=m_0}{%
    \FormulaRefAuto{x \in \{x \in A \mid P(x)\}\vdash P(x)}{8}}

  \proofstepwidestar{0=0}{\rII}
  \proofstepwidestar[1,3]{a=m_0}{\rRE{10,9}}
\end{tabproofwide}

\subparagraph{Technisches Lemma: Fixierungsmenge einer Nachfolgerstufe}

\par\Needspace{14\baselineskip}
\FormulaDefDeltaR[Fixierungsmenge der Nachfolgerstufe]{%
  \begin{aligned}[t]
    &m_0\in M\dsep f\colon M\to M\dsep x\in\mathbb{N}\dsep u\in M \vdash{}\\
    &\Phi_{x,u}\coloneq
    \{(y,d)\in \RecCore{m_0}{f}\mid y=\Succ(x)\rightarrow d=f(u)\}
  \end{aligned}
}{%
  m_0\in M\dsep f\colon M\to M\dsep x\in\mathbb{N}\dsep u\in M \vdash
  \Phi_{x,u}\coloneq
  \{(y,d)\in \RecCore{m_0}{f}\mid y=\Succ(x)\rightarrow d=f(u)\}
}{%
  \DeltaDecl{Mengen}{M\dsep x\dsep y\dsep d\dsep u\dsep m_0}%
  \DeltaDecl{Funktionen}{f}%
}

\par\Needspace{14\baselineskip}
\FormulaThmDeltaR[Fixierungsmenge der Nachfolgerstufe ist rekursionsadmissibel]{%
  \begin{aligned}[t]
    &m_0\in M\dsep f\colon M\to M\dsep x\in\mathbb{N}\dsep u\in M\dsep{}\\
    &(x,u)\in \RecCore{m_0}{f}\dsep{}\\
    &\forall c\in M\,\bigl((x,c)\in \RecCore{m_0}{f}\rightarrow c=u\bigr)\vdash{}\\
    &\AdmRec{m_0}{f}(\Phi_{x,u})
  \end{aligned}
}{%
  m_0\in M\dsep f\colon M\to M\dsep x\in\mathbb{N}\dsep u\in M\dsep
  (x,u)\in \RecCore{m_0}{f}\dsep
  \forall c\in M\,\bigl((x,c)\in \RecCore{m_0}{f}\rightarrow c=u\bigr)
  \vdash \AdmRec{m_0}{f}(\Phi_{x,u})
}{%
  \DeltaDecl{Mengen}{M\dsep x\dsep y\dsep d\dsep c\dsep u\dsep m_0}%
  \DeltaDecl{Funktionen}{f}%
}
\begin{tabproofsplitwide}
  \proofpartwideindR[Elimination aus \(\Phi_{x,u}\): Kernzugehörigkeit]{%
    \begin{aligned}[t]
      &m_0\in M\dsep f\colon M\to M\dsep x\in\mathbb{N}\dsep u\in M\dsep{}\\
      &(y,d)\in \Phi_{x,u}\vdash (y,d)\in \RecCore{m_0}{f}
    \end{aligned}
  }{%
    m_0\in M\dsep f\colon M\to M\dsep x\in\mathbb{N}\dsep u\in M\dsep
    (y,d)\in \Phi_{x,u}\vdash (y,d)\in \RecCore{m_0}{f}
  }
    \proofstepwidestar[1]{m_0\in M}{\rA}
    \proofstepwidestar[2]{f\colon M\to M}{\rA}
    \proofstepwidestar[3]{x\in\mathbb{N}}{\rA}
    \proofstepwidestar[4]{u\in M}{\rA}
    \proofstepwidestar[5]{(y,d)\in \Phi_{x,u}}{\rA}
    \proofstepwidestar[1,2,3,4]{%
      \begin{aligned}[t]
        &\Phi_{x,u}\coloneq{}\\
        &\{(y,d)\in \RecCore{m_0}{f}\mid{}\\
        &y=\Succ(x)\rightarrow d=f(u)\}
      \end{aligned}
    }{%
      \FormulaRefAuto{m_0\in M\dsep f\colon M\to M\dsep x\in\mathbb{N}\dsep u\in M \vdash \Phi_{x,u}\coloneq \{(y,d)\in \RecCore{m_0}{f}\mid y=\Succ(x)\rightarrow d=f(u)\}}{1,2,3,4}}
    \proofstepwidestar[1,2,3,4,5]{%
      \begin{aligned}[t]
        &(y,d)\in{}\\
        &\{(y,d)\in \RecCore{m_0}{f}\mid{}\\
        &y=\Succ(x)\rightarrow d=f(u)\}
      \end{aligned}
    }{\rIE{6,5}}
    \proofstepwidestar[1,2,3,4,5]{(y,d)\in \RecCore{m_0}{f}}{%
      \FormulaRefAuto{x \in \{x \in A \mid P(x)\}\vdash x\in A}{7}}
  \closeproofpartwideindR

  \proofpartwideindR[Einführung in \(\Phi_{x,u}\)]{%
    \begin{aligned}[t]
      &m_0\in M\dsep f\colon M\to M\dsep x\in\mathbb{N}\dsep u\in M\dsep{}\\
      &(z,e)\in \RecCore{m_0}{f}\dsep
      z=\Succ(x)\rightarrow e=f(u)\vdash{}\\
      &(z,e)\in \Phi_{x,u}
    \end{aligned}
  }{%
    m_0\in M\dsep f\colon M\to M\dsep x\in\mathbb{N}\dsep u\in M\dsep
    (z,e)\in \RecCore{m_0}{f}\dsep
    z=\Succ(x)\rightarrow e=f(u)\vdash (z,e)\in \Phi_{x,u}
  }
    \proofstepwidestar[1]{m_0\in M}{\rA}
    \proofstepwidestar[2]{f\colon M\to M}{\rA}
    \proofstepwidestar[3]{x\in\mathbb{N}}{\rA}
    \proofstepwidestar[4]{u\in M}{\rA}
    \proofstepwidestar[5]{(z,e)\in \RecCore{m_0}{f}}{\rA}
    \proofstepwidestar[6]{z=\Succ(x)\rightarrow e=f(u)}{\rA}
    \proofstepwidestar[5,6]{%
      \begin{aligned}[t]
        &(z,e)\in{}\\
        &\{(y,d)\in \RecCore{m_0}{f}\mid{}\\
        &y=\Succ(x)\rightarrow d=f(u)\}
      \end{aligned}
    }{%
      \FormulaRefAuto{x \in A\dsep P(x)\vdash x \in \{x \in A \mid P(x)\}}{5,6}}
    \proofstepwidestar[1,2,3,4]{%
      \begin{aligned}[t]
        &\Phi_{x,u}\coloneq{}\\
        &\{(y,d)\in \RecCore{m_0}{f}\mid{}\\
        &y=\Succ(x)\rightarrow d=f(u)\}
      \end{aligned}
    }{%
      \FormulaRefAuto{m_0\in M\dsep f\colon M\to M\dsep x\in\mathbb{N}\dsep u\in M \vdash \Phi_{x,u}\coloneq \{(y,d)\in \RecCore{m_0}{f}\mid y=\Succ(x)\rightarrow d=f(u)\}}{1,2,3,4}}
    \proofstepwidestar[1,2,3,4,5,6]{(z,e)\in \Phi_{x,u}}{%
      \rIE{8,7}}
  \closeproofpartwideindR

  \proofpartwideindR[Teilmenge des Produktraums]{%
    \begin{aligned}[t]
      &m_0\in M\dsep f\colon M\to M\dsep x\in\mathbb{N}\dsep u\in M \vdash{}\\
      &\Phi_{x,u}\subseteq \mathbb{N}\times M
    \end{aligned}
  }{%
    m_0\in M\dsep f\colon M\to M\dsep x\in\mathbb{N}\dsep u\in M \vdash
    \Phi_{x,u}\subseteq \mathbb{N}\times M
  }
    \proofstepwidestar[1]{m_0\in M}{\rA}
    \proofstepwidestar[2]{f\colon M\to M}{\rA}
    \proofstepwidestar[3]{x\in\mathbb{N}}{\rA}
    \proofstepwidestar[4]{u\in M}{\rA}
    \proofstepwidestar[1,2,3,4]{%
      \begin{aligned}[t]
        &\Phi_{x,u}\coloneq{}\\
        &\{(y,d)\in \RecCore{m_0}{f}\mid{}\\
        &y=\Succ(x)\rightarrow d=f(u)\}
      \end{aligned}
    }{%
      \FormulaRefAuto{m_0\in M\dsep f\colon M\to M\dsep x\in\mathbb{N}\dsep u\in M \vdash \Phi_{x,u}\coloneq \{(y,d)\in \RecCore{m_0}{f}\mid y=\Succ(x)\rightarrow d=f(u)\}}{1,2,3,4}}
    \proofstepwidestar{%
      \begin{aligned}[t]
        &\{(y,d)\in \RecCore{m_0}{f}\mid{}\\
        &y=\Succ(x)\rightarrow d=f(u)\}\\
        &\subseteq \RecCore{m_0}{f}
      \end{aligned}
    }{%
      \FormulaRefAuto{\{ x \in A \mid P(x) \} \subseteq A}}
    \proofstepwidestar[1,2,3,4]{%
      \Phi_{x,u}\subseteq \RecCore{m_0}{f}
    }{\rIE{5,6}}
    \proofstepwidestar[1,2]{\RecCore{m_0}{f}\subseteq \mathbb{N}\times M}{%
      \FormulaRefAuto{m_0\in M\dsep f\colon M\to M \vdash \RecCore{m_0}{f}\subseteq \mathbb{N}\times M}{1,2}}
    \proofstepwidestar[1,2,3,4]{\Phi_{x,u}\subseteq \mathbb{N}\times M}{%
      \FormulaRefAuto{A \subseteq B, B \subseteq C \vdash A \subseteq C}{7,8}}
  \closeproofpartwideindR

  \proofpartwideindR[Startpaar]{%
    \begin{aligned}[t]
      &m_0\in M\dsep f\colon M\to M\dsep x\in\mathbb{N}\dsep u\in M \vdash{}\\
      &(0,m_0)\in \Phi_{x,u}
    \end{aligned}
  }{%
    m_0\in M\dsep f\colon M\to M\dsep x\in\mathbb{N}\dsep u\in M \vdash
    (0,m_0)\in \Phi_{x,u}
  }
    \proofstepwidestar[1]{m_0\in M}{\rA}
    \proofstepwidestar[2]{f\colon M\to M}{\rA}
    \proofstepwidestar[3]{x\in\mathbb{N}}{\rA}
    \proofstepwidestar[4]{u\in M}{\rA}
    \proofstepwidestar[3]{\Succ(x)\neq0}{%
      \FormulaRefAuto{PeanoSuccNonzero}{3}}
    \proofstepwidestar[6]{0=\Succ(x)}{\rA}
    \proofstepwidestar[6]{\Succ(x)=0}{%
      \FormulaRefAuto{a=b\vdash b=a}{6}}
    \proofstepwidestar[3]{0=\Succ(x)\rightarrow \Succ(x)=0}{%
      \rRI{6,7}}
    \proofstepwidestar[3]{\neg(0=\Succ(x))}{%
      \FormulaRefAuto{P \rightarrow Q, \neg Q \vdash \neg P}{8,5}}
    \proofstepwidestar[3]{0=\Succ(x)\rightarrow m_0=f(u)}{%
      \FormulaRefAuto{\neg P \vdash P \rightarrow Q}{9}}
    \proofstepwidestar[1,2]{(0,m_0)\in \RecCore{m_0}{f}}{%
      \FormulaRefAuto{m_0\in M\dsep f\colon M\to M \vdash (0,m_0)\in \RecCore{m_0}{f}}{1,2}}
    \proofstepwidestar[1,2,3,4]{%
      (0,m_0)\in \Phi_{x,u}
    }{%
      \FormulaRefAuto{m_0\in M\dsep f\colon M\to M\dsep x\in\mathbb{N}\dsep u\in M\dsep (z,e)\in \RecCore{m_0}{f}\dsep z=\Succ(x)\rightarrow e=f(u)\vdash (z,e)\in \Phi_{x,u}}{1,2,3,4,11,10}}
  \closeproofpartwideindR

  \proofpartwideindR[Nachfolgerbarriere]{%
    \begin{aligned}[t]
      &m_0\in M\dsep f\colon M\to M\dsep x\in\mathbb{N}\dsep u\in M\dsep{}\\
      &\forall c\in M\,\bigl((x,c)\in \RecCore{m_0}{f}\rightarrow c=u\bigr)\dsep{}\\
      &y\in\mathbb{N}\dsep d\in M\dsep (y,d)\in \Phi_{x,u}\vdash{}\\
      &\Succ(y)=\Succ(x)\rightarrow f(d)=f(u)
    \end{aligned}
  }{%
    m_0\in M\dsep f\colon M\to M\dsep x\in\mathbb{N}\dsep u\in M\dsep
    \forall c\in M\,\bigl((x,c)\in \RecCore{m_0}{f}\rightarrow c=u\bigr)\dsep
    y\in\mathbb{N}\dsep d\in M\dsep (y,d)\in \Phi_{x,u}
    \vdash \Succ(y)=\Succ(x)\rightarrow f(d)=f(u)
  }
    \proofstepwidestar[1]{m_0\in M}{\rA}
    \proofstepwidestar[2]{f\colon M\to M}{\rA}
    \proofstepwidestar[3]{x\in\mathbb{N}}{\rA}
    \proofstepwidestar[4]{u\in M}{\rA}
    \proofstepwidestar[5]{%
      \begin{aligned}[t]
        &\forall c\in M\,\bigl(
        (x,c)\in \RecCore{m_0}{f}\\
        &\rightarrow c=u
        \bigr)
      \end{aligned}
    }{\rA}
    \proofstepwidestar[6]{y\in\mathbb{N}}{\rA}
    \proofstepwidestar[7]{d\in M}{\rA}
    \proofstepwidestar[8]{(y,d)\in \Phi_{x,u}}{\rA}
    \proofstepwidestar[1,2,3,4,8]{(y,d)\in \RecCore{m_0}{f}}{%
      \FormulaRefAuto{m_0\in M\dsep f\colon M\to M\dsep x\in\mathbb{N}\dsep u\in M\dsep (y,d)\in \Phi_{x,u}\vdash (y,d)\in \RecCore{m_0}{f}}{1,2,3,4,8}}
    \proofstepwidestar[10]{\Succ(y)=\Succ(x)}{\rA}
    \proofstepwidestar[10]{\Succ(x)=\Succ(y)}{%
      \FormulaRefAuto{a=b\vdash b=a}{10}}
    \proofstepwidestar[3,6,11]{x=y}{%
      \FormulaRefAuto{PeanoSuccInjective}{3,6,11}}
    \proofstepwidestar[3,6,7,10]{(x,d)=(y,d)}{%
      \FormulaRefAuto{a=c \dsep b=d\vdash (a,b)=(c,d)}{12,\rII}}
    \proofstepwidestar[1,2,3,4,6,7,8,10]{(x,d)\in \RecCore{m_0}{f}}{%
      \FormulaRefAuto{a=b \dsep b\in A \vdash a\in A}{13,9}}
    \proofstepwidestar[5,7]{(x,d)\in \RecCore{m_0}{f}\rightarrow d=u}{%
      \rRE{\rUE{5},7}}
    \proofstepwidestar[1,2,3,4,5,6,7,8,10]{d=u}{\rRE{15,14}}
    \proofstepwidestar[1,2,3,4,5,6,7,8,10]{f(d)=f(u)}{%
      \FormulaRefAuto{%
        F\colon A\to B\dsep x\in A\dsep y\in A\dsep x=y
        \vdash F(x)=F(y)%
      }{2,7,4,16}}
    \proofstepwidestar[1,2,3,4,5,6,7,8]{\Succ(y)=\Succ(x)\rightarrow f(d)=f(u)}{%
      \rRI{10,17}}
  \closeproofpartwideindR

  \proofpartwideindR[Rekursionsschritt]{%
    \begin{aligned}[t]
      &m_0\in M\dsep f\colon M\to M\dsep x\in\mathbb{N}\dsep u\in M\dsep{}\\
      &\forall c\in M\,\bigl((x,c)\in \RecCore{m_0}{f}\rightarrow c=u\bigr)\dsep{}\\
      &y\in\mathbb{N}\dsep d\in M\dsep (y,d)\in \Phi_{x,u}\vdash{}\\
      &(\Succ(y),f(d))\in \Phi_{x,u}
    \end{aligned}
  }{%
    m_0\in M\dsep f\colon M\to M\dsep x\in\mathbb{N}\dsep u\in M\dsep
    \forall c\in M\,\bigl((x,c)\in \RecCore{m_0}{f}\rightarrow c=u\bigr)\dsep
    y\in\mathbb{N}\dsep d\in M\dsep (y,d)\in \Phi_{x,u}
    \vdash (\Succ(y),f(d))\in \Phi_{x,u}
  }
    \proofstepwidestar[1]{m_0\in M}{\rA}
    \proofstepwidestar[2]{f\colon M\to M}{\rA}
    \proofstepwidestar[3]{x\in\mathbb{N}}{\rA}
    \proofstepwidestar[4]{u\in M}{\rA}
    \proofstepwidestar[5]{%
      \begin{aligned}[t]
        &\forall c\in M\,\bigl(
        (x,c)\in \RecCore{m_0}{f}\\
        &\rightarrow c=u
        \bigr)
      \end{aligned}
    }{\rA}
    \proofstepwidestar[6]{y\in\mathbb{N}}{\rA}
    \proofstepwidestar[7]{d\in M}{\rA}
    \proofstepwidestar[8]{(y,d)\in \Phi_{x,u}}{\rA}
    \proofstepwidestar[1,2,3,4,8]{(y,d)\in \RecCore{m_0}{f}}{%
      \FormulaRefAuto{m_0\in M\dsep f\colon M\to M\dsep x\in\mathbb{N}\dsep u\in M\dsep (y,d)\in \Phi_{x,u}\vdash (y,d)\in \RecCore{m_0}{f}}{1,2,3,4,8}}
    \proofstepwidestar[1,2,6,7,9]{(\Succ(y),f(d))\in \RecCore{m_0}{f}}{%
      \FormulaRefAuto{m_0\in M\dsep f\colon M\to M\dsep n\in\mathbb{N}\dsep a\in M\dsep (n,a)\in \RecCore{m_0}{f}\vdash (\Succ(n),f(a))\in \RecCore{m_0}{f}}{1,2,6,7,9}}
    \proofstepwidestar[1,2,3,4,5,6,7,8]{%
      \Succ(y)=\Succ(x)\rightarrow f(d)=f(u)
    }{%
      \FormulaRefAuto{m_0\in M\dsep f\colon M\to M\dsep x\in\mathbb{N}\dsep u\in M\dsep \forall c\in M\,\bigl((x,c)\in \RecCore{m_0}{f}\rightarrow c=u\bigr)\dsep y\in\mathbb{N}\dsep d\in M\dsep (y,d)\in \Phi_{x,u}\vdash \Succ(y)=\Succ(x)\rightarrow f(d)=f(u)}{1,2,3,4,5,6,7,8}}
    \proofstepwidestar[1,2,3,4,5,6,7,8]{%
      (\Succ(y),f(d))\in \Phi_{x,u}
    }{%
      \FormulaRefAuto{m_0\in M\dsep f\colon M\to M\dsep x\in\mathbb{N}\dsep u\in M\dsep (z,e)\in \RecCore{m_0}{f}\dsep z=\Succ(x)\rightarrow e=f(u)\vdash (z,e)\in \Phi_{x,u}}{1,2,3,4,10,11}}
  \closeproofpartwideindR

  \proofpartwide{Schluss}
    \proofstepwidestar[1]{m_0\in M}{\rA}
    \proofstepwidestar[2]{f\colon M\to M}{\rA}
    \proofstepwidestar[3]{x\in\mathbb{N}}{\rA}
    \proofstepwidestar[4]{u\in M}{\rA}
    \proofstepwidestar[5]{(x,u)\in \RecCore{m_0}{f}}{\rA}
    \proofstepwidestar[6]{%
      \begin{aligned}[t]
        &\forall c\in M\,\bigl(
        (x,c)\in \RecCore{m_0}{f}\\
        &\rightarrow c=u
        \bigr)
      \end{aligned}
    }{\rA}
    \proofstepwidestar[1,2,3,4]{(0,m_0)\in \Phi_{x,u}}{%
      \FormulaRefAuto{m_0\in M\dsep f\colon M\to M\dsep x\in\mathbb{N}\dsep u\in M \vdash (0,m_0)\in \Phi_{x,u}}{1,2,3,4}}
    \proofstepwidestar[1,2,3,4]{\Phi_{x,u}\subseteq \mathbb{N}\times M}{%
      \FormulaRefAuto{m_0\in M\dsep f\colon M\to M\dsep x\in\mathbb{N}\dsep u\in M \vdash \Phi_{x,u}\subseteq \mathbb{N}\times M}{1,2,3,4}}
    \proofstepwidestar[1,2,3,4,6]{\forall y\in\mathbb N\,\forall d\in M\,\bigl((y,d)\in \Phi_{x,u}\rightarrow (\Succ(y),f(d))\in \Phi_{x,u}\bigr)}{%
      \FormulaRefAuto{m_0\in M\dsep f\colon M\to M\dsep x\in\mathbb{N}\dsep u\in M\dsep \forall c\in M\,\bigl((x,c)\in \RecCore{m_0}{f}\rightarrow c=u\bigr)\dsep y\in\mathbb{N}\dsep d\in M\dsep (y,d)\in \Phi_{x,u}\vdash (\Succ(y),f(d))\in \Phi_{x,u}}{1,2,3,4,6}}
    \proofstepwidestar[1,2,3,4,5,6]{\AdmRec{m_0}{f}(\Phi_{x,u})}{\FormulaRefAuto{\AdmRec{m_0}{f}(H)}{8,7,9}}
  \closeproofpartwide
\end{tabproofsplitwide}

\par\Needspace{14\baselineskip}
\FormulaThmDeltaR[Die Nachfolgerstufe des Rekursionskerns ist durch die Vorstufe festgelegt]{%
  \begin{aligned}[t]
    &m_0\in M\dsep f\colon M\to M\dsep n\in\mathbb{N}\dsep a\in M\dsep{}\\
    &(n,a)\in \RecCore{m_0}{f}\dsep{}\\
    &\forall c\in M\,\bigl((n,c)\in \RecCore{m_0}{f}\rightarrow c=a\bigr)\dsep{}\\
    &(\Succ(n),b)\in \RecCore{m_0}{f}\vdash b=f(a)
  \end{aligned}
}{%
  m_0\in M\dsep f\colon M\to M\dsep n\in\mathbb{N}\dsep a\in M\dsep
  (n,a)\in \RecCore{m_0}{f}\dsep
  \forall c\in M\,\bigl((n,c)\in \RecCore{m_0}{f}\rightarrow c=a\bigr)\dsep
  (\Succ(n),b)\in \RecCore{m_0}{f}\vdash b=f(a)
}{%
  \DeltaDecl{Mengen}{M\dsep n\dsep a\dsep b\dsep c\dsep m_0}%
  \DeltaDecl{Funktionen}{f}%
}
\begin{tabproofwide}
  \proofstepwidestar[1]{m_0\in M}{\rA}
  \proofstepwidestar[2]{f\colon M\to M}{\rA}
  \proofstepwidestar[3]{n\in\mathbb{N}}{\rA}
  \proofstepwidestar[4]{a\in M}{\rA}
  \proofstepwidestar[5]{(n,a)\in \RecCore{m_0}{f}}{\rA}
  \proofstepwidestar[6]{%
    \forall c\in M\,\bigl((n,c)\in \RecCore{m_0}{f}\rightarrow c=a\bigr)
  }{\rA}
  \proofstepwidestar[7]{(\Succ(n),b)\in \RecCore{m_0}{f}}{\rA}

  \proofstepwidestar[1,2,3,4,5,6]{\AdmRec{m_0}{f}(\Phi_{n,a})}{%
    \FormulaRefAuto{m_0\in M\dsep f\colon M\to M\dsep x\in\mathbb{N}\dsep u\in M\dsep (x,u)\in \RecCore{m_0}{f}\dsep \forall c\in M\,\bigl((x,c)\in \RecCore{m_0}{f}\rightarrow c=u\bigr)\vdash \AdmRec{m_0}{f}(\Phi_{x,u})}{1,2,3,4,5,6}}

  \proofstepwidestar[1,2]{\RecCore{m_0}{f}\subseteq \Phi_{n,a}}{%
    \FormulaRefAuto{m_0\in M\dsep f\colon M\to M\dsep \AdmRec{m_0}{f}(H)\vdash \RecCore{m_0}{f}\subseteq H}{1,2,8}}

  \proofstepwidestar[1,2,3,4,5,6,7]{(\Succ(n),b)\in \Phi_{n,a}}{%
    \FormulaRefAuto{A\subseteq B\dsep x\in A\vdash x\in B}{9,7}}

  \proofstepwidestar[1,2,3,4]{%
    \begin{aligned}[t]
      &\Phi_{n,a}\coloneq{}\\
      &\{(y,d)\in \RecCore{m_0}{f}\mid{}\\
      &\qquad y=\Succ(n)\rightarrow d=f(a)\}
    \end{aligned}
  }{%
    \FormulaRefAuto{m_0\in M\dsep f\colon M\to M\dsep x\in\mathbb{N}\dsep u\in M \vdash \Phi_{x,u}\coloneq \{(y,d)\in \RecCore{m_0}{f}\mid y=\Succ(x)\rightarrow d=f(u)\}}{1,2,3,4}}

  \proofstepwidestar[1,2,3,4,5,6,7]{%
    \begin{aligned}[t]
      &(\Succ(n),b)\in{}\\
      &\{(y,d)\in \RecCore{m_0}{f}\mid{}\\
      &\qquad y=\Succ(n)\rightarrow d=f(a)\}
    \end{aligned}
  }{%
    \rIE{11,10}}

  \proofstepwidestar[1,2,3,4,5,6,7]{%
    \Succ(n)=\Succ(n)\rightarrow b=f(a)
  }{%
    \FormulaRefAuto{x \in \{x \in A \mid P(x)\}\vdash P(x)}{12}}

  \proofstepwidestar{\Succ(n)=\Succ(n)}{\rII}

  \proofstepwidestar[1,2,3,4,5,6,7]{b=f(a)}{\rRE{14,13}}
\end{tabproofwide}

\par\Needspace{14\baselineskip}
\FormulaThmDeltaK[Fixierungslemma für Rekursionskernstufen]{%
  \begin{aligned}[t]
    &m_0\in M\dsep f\colon M\to M\dsep x\in\mathbb{N}\dsep u\in M\dsep{}\\
    &(x,u)\in \RecCore{m_0}{f}\dsep
    \forall c\in M\,\bigl((x,c)\in \RecCore{m_0}{f}\rightarrow c=u\bigr)
    \vdash{}\\
    &\bigl((0,a)\in \RecCore{m_0}{f}\rightarrow a=m_0\bigr)\land{}\\
    &\bigl((\Succ(x),b)\in \RecCore{m_0}{f}\rightarrow b=f(u)\bigr)
  \end{aligned}
}{RecCoreStageBarrier}{%
  \DeltaDecl{Mengen}{M\dsep x\dsep u\dsep a\dsep b\dsep c\dsep m_0}%
  \DeltaDecl{Funktionen}{f}%
}
\begin{tabproofsplitwide}
\proofpartwide{Fixierung der Nullstufe}
\proofstepwidestar[1]{m_0\in M}{\rA}
  \proofstepwidestar[2]{f\colon M\to M}{\rA}
  \proofstepwidestar[3]{x\in\mathbb{N}}{\rA}
  \proofstepwidestar[4]{u\in M}{\rA}
  \proofstepwidestar[5]{(x,u)\in \RecCore{m_0}{f}}{\rA}
  \proofstepwidestar[6]{%
    \forall c\in M\,\bigl((x,c)\in \RecCore{m_0}{f}\rightarrow c=u\bigr)
  }{\rA}

  \proofstepwidestar[7]{(0,a)\in \RecCore{m_0}{f}}{\rA}
  \proofstepwidestar[1,2,7]{a=m_0}{%
    \FormulaRefAuto{m_0\in M\dsep f\colon M\to M\dsep
    (0,a)\in \RecCore{m_0}{f}\vdash a=m_0}{1,2,7}}
  \proofstepwidestar[1,2]{(0,a)\in \RecCore{m_0}{f}\rightarrow a=m_0}{%
    \rRI{7,8}}


\closeproofpartwide

\proofpartwide{Fixierung der Nachfolgerstufe}
  \setcounter{proofstepnr}{9}% Fortlaufende Schritte dieses Beweises.
\proofstepwidestar[10]{(\Succ(x),b)\in \RecCore{m_0}{f}}{\rA}
  \proofstepwidestar[1,2,3,4,5,6,10]{b=f(u)}{%
    \FormulaRefAuto{m_0\in M\dsep f\colon M\to M\dsep n\in\mathbb{N}\dsep a\in M\dsep
    (n,a)\in \RecCore{m_0}{f}\dsep
    \forall c\in M\,\bigl((n,c)\in \RecCore{m_0}{f}\rightarrow c=a\bigr)\dsep
    (\Succ(n),b)\in \RecCore{m_0}{f}\vdash b=f(a)}{1,2,3,4,5,6,10}}
  \proofstepwidestar[1,2,3,4,5,6]{%
    (\Succ(x),b)\in \RecCore{m_0}{f}\rightarrow b=f(u)
  }{\rRI{10,11}}


\closeproofpartwide

\proofpartwide{Zusammenführung}
  \setcounter{proofstepnr}{12}% Fortlaufende Schritte dieses Beweises.
\proofstepwidestar[1,2,3,4,5,6]{%
    \bigl((0,a)\in \RecCore{m_0}{f}\rightarrow a=m_0\bigr)\land
    \bigl((\Succ(x),b)\in \RecCore{m_0}{f}\rightarrow b=f(u)\bigr)
  }{\rAI{9,12}}

\closeproofpartwide
\end{tabproofsplitwide}

\subsubsection{Das Stufenlemma}

Dieser Block ist der eigentliche Punkt der Konstruktion. Statt Existenz und
Eindeutigkeit jeder Stufe getrennt zu beweisen, wird direkt über das Prädikat
\[
  U(n)\;:\Longleftrightarrow\;
  \exists!a\in M\,\bigl((n,a)\in\RecCore{m_0}{f}\bigr)
\]
induktiv argumentiert.

\par\Needspace{14\baselineskip}
\FormulaThmDeltaK[Stufenlemma des Rekursionskerns]{%
  m_0\in M\dsep f\colon M\to M\dsep n\in\mathbb{N}
  \vdash \exists! a\in M\,\bigl((n,a)\in \RecCore{m_0}{f}\bigr)
}{RecCoreStageUnique}{%
  \DeltaDecl{Mengen}{M\dsep n\dsep x\dsep a\dsep b\dsep c\dsep u\dsep m_0}%
  \DeltaDecl{Funktionen}{f}%
}
\begin{tabproofsplitwide}
  \proofpartwideindR[Induktionsanfang]{%
    m_0\in M\dsep f\colon M\to M
    \vdash \exists! a\in M\,\bigl((0,a)\in \RecCore{m_0}{f}\bigr)
  }{%
    m_0\in M\dsep f\colon M\to M
    \vdash \exists! a\in M\,\bigl((0,a)\in \RecCore{m_0}{f}\bigr)
  }
    \proofstepwidestar[1]{m_0\in M}{\rA}
    \proofstepwidestar[2]{f\colon M\to M}{\rA}
    \proofstepwidestar[1,2]{(0,m_0)\in \RecCore{m_0}{f}}{%
      \FormulaRefAuto{m_0\in M\dsep f\colon M\to M \vdash (0,m_0)\in \RecCore{m_0}{f}}{1,2}}
    \proofstepwidestar[1,2]{%
      \exists a\in M\,\bigl((0,a)\in \RecCore{m_0}{f}\bigr)
    }{\rEI{\rAI{1,3}}}

    \proofstepwidestar[5]{u\in M}{\rA}
    \proofstepwidestar[6]{c\in M}{\rA}
    \proofstepwidestar[7]{(0,u)\in \RecCore{m_0}{f}}{\rA}
    \proofstepwidestar[8]{(0,c)\in \RecCore{m_0}{f}}{\rA}
    \proofstepwidestar[1,2,7]{u=m_0}{%
      \FormulaRefAuto{m_0\in M\dsep f\colon M\to M\dsep
      (0,a)\in \RecCore{m_0}{f}\vdash a=m_0}{1,2,7}}
    \proofstepwidestar[1,2,8]{c=m_0}{%
      \FormulaRefAuto{m_0\in M\dsep f\colon M\to M\dsep
      (0,a)\in \RecCore{m_0}{f}\vdash a=m_0}{1,2,8}}
    \proofstepwidestar[1,2,7,8]{u=c}{%
      \FormulaRefAuto{a = b,\, c = b \vdash a = c}{9,10}}

    \proofstepwidestar[1,2]{%
      \exists! a\in M\,\bigl((0,a)\in \RecCore{m_0}{f}\bigr)
    }{\UEI{4,5,6,7,8,11}}
  \closeproofpartwideindR

  \proofpartwideindR[Induktionsschritt]{%
    \begin{aligned}[t]
      &m_0\in M\dsep f\colon M\to M\dsep x\in\mathbb{N}\dsep{}\\
      &\exists! a\in M\,\bigl((x,a)\in \RecCore{m_0}{f}\bigr)
      \vdash{}\\
      &\exists! b\in M\,\bigl((\Succ(x),b)\in \RecCore{m_0}{f}\bigr)
    \end{aligned}
  }{%
    m_0\in M\dsep f\colon M\to M\dsep x\in\mathbb{N}\dsep
    \exists! a\in M\,\bigl((x,a)\in \RecCore{m_0}{f}\bigr)
    \vdash
    \exists! b\in M\,\bigl((\Succ(x),b)\in \RecCore{m_0}{f}\bigr)
  }
    \proofstepwidestar[1]{m_0\in M}{\rA}
    \proofstepwidestar[2]{f\colon M\to M}{\rA}
    \proofstepwidestar[3]{x\in\mathbb{N}}{\rA}
    \proofstepwidestar[4]{%
      \exists! a\in M\,\bigl((x,a)\in \RecCore{m_0}{f}\bigr)
    }{\rA}

    \proofstepwidestar[5]{u\in M}{\rA}
    \proofstepwidestar[6]{(x,u)\in \RecCore{m_0}{f}}{\rA}
    \proofstepwidestar[1,2,3,5,6]{(\Succ(x),f(u))\in \RecCore{m_0}{f}}{%
      \FormulaRefAuto{m_0\in M\dsep f\colon M\to M\dsep n\in\mathbb{N}\dsep a\in M\dsep
      (n,a)\in \RecCore{m_0}{f}
      \vdash (\Succ(n),f(a))\in \RecCore{m_0}{f}}{1,2,3,5,6}}
    \proofstepwidestar[2,5]{f ( u ) \in M}{%
      \FormulaRefAuto{F\colon A\to B\dsep x\in A\vdash F(x)\in B}{2,5}}
    \proofstepwidestar[1,2,3,5,6]{%
      \exists b\in M\,\bigl((\Succ(x),b)\in \RecCore{m_0}{f}\bigr)
    }{\rEI{\rAI{%
      8,7}}}

    \proofstepwidestar[4,5,6]{%
      \forall c\in M\,\bigl((x,c)\in \RecCore{m_0}{f}\rightarrow c=u\bigr)
    }{\UEE{4,5,6}}

    \proofstepwidestar[11]{b\in M}{\rA}
    \proofstepwidestar[12]{c\in M}{\rA}
    \proofstepwidestar[13]{(\Succ(x),b)\in \RecCore{m_0}{f}}{\rA}
    \proofstepwidestar[14]{(\Succ(x),c)\in \RecCore{m_0}{f}}{\rA}
    \proofstepwidestar[1,2,3,5,6,10,13]{b=f(u)}{%
      \FormulaRefAuto{m_0\in M\dsep f\colon M\to M\dsep n\in\mathbb{N}\dsep a\in M\dsep
      (n,a)\in \RecCore{m_0}{f}\dsep
      \forall c\in M\,\bigl((n,c)\in \RecCore{m_0}{f}\rightarrow c=a\bigr)\dsep
      (\Succ(n),b)\in \RecCore{m_0}{f}\vdash b=f(a)}{1,2,3,5,6,10,13}}
    \proofstepwidestar[1,2,3,5,6,10,14]{c=f(u)}{%
      \FormulaRefAuto{m_0\in M\dsep f\colon M\to M\dsep n\in\mathbb{N}\dsep a\in M\dsep
      (n,a)\in \RecCore{m_0}{f}\dsep
      \forall c\in M\,\bigl((n,c)\in \RecCore{m_0}{f}\rightarrow c=a\bigr)\dsep
      (\Succ(n),b)\in \RecCore{m_0}{f}\vdash b=f(a)}{1,2,3,5,6,10,14}}
    \proofstepwidestar[1,2,3,5,6,10,13,14]{b=c}{%
      \FormulaRefAuto{a = b,\, c = b \vdash a = c}{15,16}}

    \proofstepwidestar[1,2,3,4,5,6]{%
      \exists! b\in M\,\bigl((\Succ(x),b)\in \RecCore{m_0}{f}\bigr)
    }{\UEI{9,11,12,13,14,17}}
    \proofstepwidestar[1,2,3,4]{%
      \exists! b\in M\,\bigl((\Succ(x),b)\in \RecCore{m_0}{f}\bigr)
    }{\UEE{4,5,6,18}}
  \closeproofpartwideindR

  \proofpartwide{Schluss}
    \proofstepwidestar[1]{m_0\in M}{\rA}
    \proofstepwidestar[2]{f\colon M\to M}{\rA}
    \proofstepwidestar[3]{n\in\mathbb{N}}{\rA}

    \proofstepwidestar[1,2]{\exists! a\in M\,\bigl((0,a)\in \RecCore{m_0}{f}\bigr)}{%
      \FormulaRefAuto{m_0\in M\dsep f\colon M\to M
          \vdash \exists! a\in M\,\bigl((0,a)\in \RecCore{m_0}{f}\bigr)}{1,2}}
    \proofstepwidestar[1,2]{\forall x\in\mathbb N\,\bigl((\exists! a\in M\,\bigl((x,a)\in \RecCore{m_0}{f}\bigr))\rightarrow(\exists! b\in M\,\bigl((\Succ(x),b)\in \RecCore{m_0}{f}\bigr))\bigr)}{%
      \FormulaRefAuto{m_0\in M\dsep f\colon M\to M\dsep x\in\mathbb{N}\dsep
          \exists! a\in M\,\bigl((x,a)\in \RecCore{m_0}{f}\bigr)
          \vdash
          \exists! b\in M\,\bigl((\Succ(x),b)\in \RecCore{m_0}{f}\bigr)}{1,2}}
    \proofstepwidestar[1,2,3]{%
      \exists! a\in M\,\bigl((n,a)\in \RecCore{m_0}{f}\bigr)
    }{\FormulaRefAuto{PeanoInductionRule}{4,5,3}}
  \closeproofpartwide
\end{tabproofsplitwide}

\iffalse

\paragraph{Existenz}

\par\Needspace{14\baselineskip}
\FormulaThmDelta[Existenz eines Werts in jeder Stufe]{%
  m_0\in M\dsep f\colon M\to M\dsep n\in\mathbb{N}
  \vdash \exists a\in M\,\bigl((n,a)\in \RecCore{m_0}{f}\bigr)
}{%
  \DeltaDecl{Mengen}{M\dsep n\dsep x\dsep a\dsep b\dsep m_0}%
  \DeltaDecl{Funktionen}{f}%
}
\begin{tabproofsplitwide}
  \proofpartwideindR[Induktionsanfang]{%
    m_0\in M\dsep f\colon M\to M
    \vdash \exists a\in M\,\bigl((0,a)\in \RecCore{m_0}{f}\bigr)
  }{%
    m_0\in M\dsep f\colon M\to M
    \vdash \exists a\in M\,\bigl((0,a)\in \RecCore{m_0}{f}\bigr)
  }
    \proofstepwidestar[1]{m_0\in M}{\rA}
    \proofstepwidestar[2]{f\colon M\to M}{\rA}
    \proofstepwidestar[1,2]{(0,m_0)\in \RecCore{m_0}{f}}{%
      \FormulaRefAuto{m_0\in M\dsep f\colon M\to M \vdash (0,m_0)\in \RecCore{m_0}{f}}{1,2}}
    \proofstepwidestar[1,2]{%
      \exists a\in M\,\bigl((0,a)\in \RecCore{m_0}{f}\bigr)
    }{%
      \rEI{\rAI{1,3}}}
  \closeproofpartwideindR

    \proofpartwideindR[Induktionsschritt]{%
      \begin{aligned}[t]
        &m_0\in M\dsep f\colon M\to M\dsep{}\\
        &x\in\mathbb{N}\dsep
          \exists a\in M\,\bigl((x,a)\in \RecCore{m_0}{f}\bigr)\vdash{}\\
        &\exists b\in M\,\bigl((\Succ(x),b)\in \RecCore{m_0}{f}\bigr)
      \end{aligned}
    }{%
      m_0\in M\dsep f\colon M\to M\dsep
      x\in\mathbb{N}\dsep
      \exists a\in M\,\bigl((x,a)\in \RecCore{m_0}{f}\bigr)\vdash
      \exists b\in M\,\bigl((\Succ(x),b)\in \RecCore{m_0}{f}\bigr)
    }
    \proofstepwidestar[1]{m_0\in M}{\rA}
    \proofstepwidestar[2]{f\colon M\to M}{\rA}
    \proofstepwidestar[3]{x\in\mathbb{N}}{\rA}
    \proofstepwidestar[4]{%
      \exists a\in M\,\bigl((x,a)\in \RecCore{m_0}{f}\bigr)
    }{\rA}

    \proofstepwidestar[5]{a\in M}{\rA}
    \proofstepwidestar[6]{(x,a)\in \RecCore{m_0}{f}}{\rA}

    \proofstepwidestar[2,5]{f(a)\in M}{%
      \FormulaRefAuto{F\colon A\to B\dsep x\in A\vdash F(x)\in B}{2,5}}

    \proofstepwidestar[1,2,3,5,6]{%
      (\Succ(x),f(a))\in \RecCore{m_0}{f}
    }{%
      \FormulaRefAuto{m_0\in M\dsep f\colon M\to M\dsep n\in\mathbb{N}\dsep a\in M\dsep (n,a)\in \RecCore{m_0}{f}\vdash (\Succ(n),f(a))\in \RecCore{m_0}{f}}{1,2,3,5,6}}

    \proofstepwidestar[1,2,3,5,6]{%
      \exists b\in M\,\bigl((\Succ(x),b)\in \RecCore{m_0}{f}\bigr)
    }{%
      \rEI{\rAI{7,8}}}

    \proofstepwidestar[1,2,3,4]{%
      \exists b\in M\,\bigl((\Succ(x),b)\in \RecCore{m_0}{f}\bigr)
    }{%
      \rEE{4,5,6,9}}
  \closeproofpartwideindR

  \proofpartwide{Schluss}
    \proofstepwidestar[1]{m_0\in M}{\rA}
    \proofstepwidestar[2]{f\colon M\to M}{\rA}
    \proofstepwidestar[3]{n\in\mathbb{N}}{\rA}

    \proofstepwidestar[1,2]{\exists a\in M\,\bigl((0,a)\in \RecCore{m_0}{f}\bigr)}{%
      \FormulaRefAuto{m_0\in M\dsep f\colon M\to M \vdash \exists a\in M\,\bigl((0,a)\in \RecCore{m_0}{f}\bigr)}{1,2}}
    \proofstepwidestar[1,2]{\forall x\in\mathbb N\,\bigl((\exists a\in M\,\bigl((x,a)\in \RecCore{m_0}{f}\bigr))\rightarrow(\exists b\in M\,\bigl((\Succ(x),b)\in \RecCore{m_0}{f}\bigr))\bigr)}{%
      \FormulaRefAuto{m_0\in M\dsep f\colon M\to M\dsep
      x\in\mathbb{N}\dsep
      \exists a\in M\,\bigl((x,a)\in \RecCore{m_0}{f}\bigr)\vdash
      \exists b\in M\,\bigl((\Succ(x),b)\in \RecCore{m_0}{f}\bigr)}{1,2}}
    \proofstepwidestar[1,2,3]{%
      \exists a\in M\,\bigl((n,a)\in \RecCore{m_0}{f}\bigr)
    }{\FormulaRefAuto{PeanoInductionRule}{4,5,3}}
  \closeproofpartwide
\end{tabproofsplitwide}

\paragraph{Eindeutigkeit}

\par\Needspace{14\baselineskip}
\FormulaThmDeltaR[Eindeutigkeit in der Nullstufe]{%
  \begin{aligned}[t]
    &m_0\in M\dsep f\colon M\to M\dsep{}\\
    &(0,u)\in \RecCore{m_0}{f}\dsep
      (0,c)\in \RecCore{m_0}{f}\vdash{}\\
    &u=c
  \end{aligned}
}{%
  m_0\in M\dsep f\colon M\to M\dsep
  (0,u)\in \RecCore{m_0}{f}\dsep
  (0,c)\in \RecCore{m_0}{f}\vdash u=c
}{%
  \DeltaDecl{Mengen}{M\dsep u\dsep c\dsep m_0}%
  \DeltaDecl{Funktionen}{f}%
}
\begin{tabproof}
  \proofstep{1}{m_0\in M}{\rA}
  \proofstep{2}{f\colon M\to M}{\rA}
  \proofstep{3}{(0,u)\in \RecCore{m_0}{f}}{\rA}
  \proofstep{4}{(0,c)\in \RecCore{m_0}{f}}{\rA}

  \proofstep{1,2,3}{u=m_0}{%
    \FormulaRefAuto{m_0\in M\dsep f\colon M\to M\dsep (0,a)\in \RecCore{m_0}{f}\vdash a=m_0}{1,2,3}}

  \proofstep{1,2,4}{c=m_0}{%
    \FormulaRefAuto{m_0\in M\dsep f\colon M\to M\dsep (0,a)\in \RecCore{m_0}{f}\vdash a=m_0}{1,2,4}}

  \proofstep{1,2,3,4}{u=c}{%
    \FormulaRefAuto{a = b,\, c = b \vdash a = c}{5,6}}
\end{tabproof}


\par\Needspace{14\baselineskip}
\FormulaThmDeltaR[Eindeutigkeit in der Nachfolgerstufe bei eindeutiger Vorstufe]{%
  \begin{aligned}[t]
    &m_0\in M\dsep f\colon M\to M\dsep x\in\mathbb{N}\dsep{}\\
    &\forall u\in M\,\bigl((x,u)\in \RecCore{m_0}{f}\rightarrow{}\\
    &\qquad \forall c\in M\,\bigl((x,c)\in \RecCore{m_0}{f}\rightarrow u=c\bigr)\bigr)\dsep{}\\
    &(\Succ(x),d)\in \RecCore{m_0}{f}\dsep
      (\Succ(x),v)\in \RecCore{m_0}{f}\vdash{}\\
    &d=v
  \end{aligned}
}{%
  m_0\in M\dsep f\colon M\to M\dsep x\in\mathbb{N}\dsep
  \forall u\in M\,\bigl((x,u)\in \RecCore{m_0}{f}\rightarrow
  \forall c\in M\,\bigl((x,c)\in \RecCore{m_0}{f}\rightarrow u=c\bigr)\bigr)\dsep
  (\Succ(x),d)\in \RecCore{m_0}{f}\dsep
  (\Succ(x),v)\in \RecCore{m_0}{f}\vdash d=v
}{%
  \DeltaDecl{Mengen}{M\dsep x\dsep d\dsep v\dsep u\dsep c\dsep m_0}%
  \DeltaDecl{Funktionen}{f}%
}
\begin{tabproof}
  \proofstep{1}{m_0\in M}{\rA}
  \proofstep{2}{f\colon M\to M}{\rA}
  \proofstep{3}{x\in\mathbb{N}}{\rA}
  \proofstep{4}{%
    \begin{aligned}[t]
      &\forall u\in M\,\bigl(
      (x,u)\in \RecCore{m_0}{f}\rightarrow{}\\
      &\forall c\in M\,\bigl(
      (x,c)\in \RecCore{m_0}{f}\rightarrow u=c
      \bigr)\bigr)
    \end{aligned}
  }{\rA}
  \proofstep{5}{(\Succ(x),d)\in \RecCore{m_0}{f}}{\rA}
  \proofstep{6}{(\Succ(x),v)\in \RecCore{m_0}{f}}{\rA}

  \proofstep{1,2,3}{%
    \exists u\in M\,\bigl((x,u)\in \RecCore{m_0}{f}\bigr)
  }{%
    \FormulaRefAuto{m_0\in M\dsep f\colon M\to M\dsep n\in\mathbb{N}\vdash \exists a\in M\,\bigl((n,a)\in \RecCore{m_0}{f}\bigr)}{1,2,3}}

  \proofstep{8}{u\in M}{\rA}
  \proofstep{9}{(x,u)\in \RecCore{m_0}{f}}{\rA}

  \proofstep{4,8}{%
    \begin{aligned}[t]
      &(x,u)\in \RecCore{m_0}{f}\rightarrow{}\\
      &\forall c\in M\,\bigl(
      (x,c)\in \RecCore{m_0}{f}\rightarrow u=c
      \bigr)
    \end{aligned}
  }{%
    \rRE{\rUE{4},8}}

  \proofstep{4,8,9}{%
    \begin{aligned}[t]
      &\forall c\in M\,\bigl(
      (x,c)\in \RecCore{m_0}{f}\rightarrow{}\\
      &u=c
      \bigr)
    \end{aligned}
  }{%
    \rRE{9,10}}

  \proofstep{1,2,3,8,9,11,5}{d=f(u)}{%
    \FormulaRefAuto{m_0\in M\dsep f\colon M\to M\dsep n\in\mathbb{N}\dsep a\in M\dsep (n,a)\in \RecCore{m_0}{f}\dsep \forall c\in M\,\bigl((n,c)\in \RecCore{m_0}{f}\rightarrow c=a\bigr)\dsep (\Succ(n),b)\in \RecCore{m_0}{f}\vdash b=f(a)}{1,2,3,8,9,11,5}}

  \proofstep{1,2,3,8,9,11,6}{v=f(u)}{%
    \FormulaRefAuto{m_0\in M\dsep f\colon M\to M\dsep n\in\mathbb{N}\dsep a\in M\dsep (n,a)\in \RecCore{m_0}{f}\dsep \forall c\in M\,\bigl((n,c)\in \RecCore{m_0}{f}\rightarrow c=a\bigr)\dsep (\Succ(n),b)\in \RecCore{m_0}{f}\vdash b=f(a)}{1,2,3,8,9,11,6}}

  \proofstep{1,2,3,4,5,6,8,9}{d=v}{%
    \FormulaRefAuto{a = b,\, c = b \vdash a = c}{12,13}}

  \proofstep{1,2,3,4,5,6}{d=v}{%
    \rEE{7,8,9,14}}
\end{tabproof}

\par\Needspace{14\baselineskip}
\FormulaThmDeltaR[Eindeutigkeit eines Werts in jeder Stufe]{%
  \begin{aligned}[t]
    &m_0\in M\dsep f\colon M\to M\dsep n\in\mathbb{N}\dsep{}\\
    &a\in M\dsep b\in M\dsep{}\\
    &(n,a)\in \RecCore{m_0}{f}\dsep (n,b)\in \RecCore{m_0}{f}\vdash a=b
  \end{aligned}
}{%
  m_0\in M\dsep f\colon M\to M\dsep n\in\mathbb{N}\dsep
  a\in M\dsep b\in M\dsep
  (n,a)\in \RecCore{m_0}{f}\dsep (n,b)\in \RecCore{m_0}{f}\vdash a=b
}{%
  \DeltaDecl{Mengen}{M\dsep n\dsep x\dsep a\dsep b\dsep c\dsep d\dsep u\dsep m_0}%
  \DeltaDecl{Funktionen}{f}%
}
\begin{tabproofsplitwide}
  \proofpartwideindR[Induktionsanfang]{%
    \begin{aligned}[t]
      &m_0\in M\dsep f\colon M\to M \vdash{}\\
      &\forall u\in M\,\bigl((0,u)\in \RecCore{m_0}{f}\rightarrow{}\\
      &\qquad \forall c\in M\,\bigl((0,c)\in \RecCore{m_0}{f}\rightarrow u=c\bigr)\bigr)
    \end{aligned}
  }{%
    m_0\in M\dsep f\colon M\to M \vdash
    \forall u\in M\,\bigl((0,u)\in \RecCore{m_0}{f}\rightarrow
    \forall c\in M\,\bigl((0,c)\in \RecCore{m_0}{f}\rightarrow u=c\bigr)\bigr)
  }
    \proofstepwidestar[1]{m_0\in M}{\rA}
    \proofstepwidestar[2]{f\colon M\to M}{\rA}
    \proofstepwidestar[3]{u\in M}{\rA}
    \proofstepwidestar[4]{(0,u)\in \RecCore{m_0}{f}}{\rA}
    \proofstepwidestar[5]{c\in M}{\rA}
    \proofstepwidestar[6]{(0,c)\in \RecCore{m_0}{f}}{\rA}

    \proofstepwidestar[1,2,4,6]{u=c}{%
      \FormulaRefAuto{m_0\in M\dsep f\colon M\to M\dsep (0,u)\in \RecCore{m_0}{f}\dsep (0,c)\in \RecCore{m_0}{f}\vdash u=c}{1,2,4,6}}

    \proofstepwidestar[1,2,4,5]{%
      (0,c)\in \RecCore{m_0}{f}\rightarrow u=c
    }{%
      \rRI{6,7}}

    \proofstepwidestar[1,2,4]{%
      \forall c\in M\,\bigl((0,c)\in \RecCore{m_0}{f}\rightarrow u=c\bigr)
    }{%
      \rUI{\rRI{5,8}}}

    \proofstepwidestar[1,2,3]{%
      \begin{aligned}[t]
        &(0,u)\in \RecCore{m_0}{f}\rightarrow{}\\
        &\forall c\in M\,\bigl((0,c)\in \RecCore{m_0}{f}\rightarrow u=c\bigr)
      \end{aligned}
    }{%
      \rRI{4,9}}

    \proofstepwidestar[1,2]{%
      \begin{aligned}[t]
        &\forall u\in M\,\bigl((0,u)\in \RecCore{m_0}{f}\rightarrow{}\\
        &\qquad \forall c\in M\,\bigl((0,c)\in \RecCore{m_0}{f}\rightarrow u=c\bigr)\bigr)
      \end{aligned}
    }{%
      \rUI{\rRI{3,10}}}
  \closeproofpartwideindR

  \proofpartwideindR[Induktionsschritt]{%
  \begin{aligned}[t]
    &m_0\in M\dsep f\colon M\to M\dsep x\in\mathbb{N}\dsep{}\\
    &\forall u\in M\,\bigl((x,u)\in \RecCore{m_0}{f}\rightarrow{}\\
    &\qquad \forall c\in M\,\bigl((x,c)\in \RecCore{m_0}{f}\rightarrow u=c\bigr)\bigr)\vdash{}\\
    &\forall a\in M\,\bigl((\Succ(x),a)\in \RecCore{m_0}{f}\rightarrow{}\\
    &\qquad \forall c\in M\,\bigl((\Succ(x),c)\in \RecCore{m_0}{f}\rightarrow a=c\bigr)\bigr)
  \end{aligned}
}{%
  m_0\in M\dsep f\colon M\to M\dsep x\in\mathbb{N}\dsep
  \forall u\in M\,\bigl((x,u)\in \RecCore{m_0}{f}\rightarrow
  \forall c\in M\,\bigl((x,c)\in \RecCore{m_0}{f}\rightarrow u=c\bigr)\bigr)\vdash
  \forall a\in M\,\bigl((\Succ(x),a)\in \RecCore{m_0}{f}\rightarrow
  \forall c\in M\,\bigl((\Succ(x),c)\in \RecCore{m_0}{f}\rightarrow a=c\bigr)\bigr)
}
    \proofstepwidestar[1]{m_0\in M}{\rA}
    \proofstepwidestar[2]{f\colon M\to M}{\rA}
    \proofstepwidestar[3]{x\in\mathbb{N}}{\rA}
    \proofstepwidestar[4]{%
      \begin{aligned}[t]
        &\forall u\in M\,\bigl((x,u)\in \RecCore{m_0}{f}\rightarrow{}\\
        &\qquad \forall c\in M\,\bigl((x,c)\in \RecCore{m_0}{f}\rightarrow u=c\bigr)\bigr)
      \end{aligned}
    }{\rA}

    \proofstepwidestar[5]{a\in M}{\rA}
    \proofstepwidestar[6]{(\Succ(x),a)\in \RecCore{m_0}{f}}{\rA}
    \proofstepwidestar[7]{c\in M}{\rA}
    \proofstepwidestar[8]{(\Succ(x),c)\in \RecCore{m_0}{f}}{\rA}

    \proofstepwidestar[1,2,3,4,6,8]{a=c}{%
      \FormulaRefAuto{%
        m_0\in M\dsep f\colon M\to M\dsep x\in\mathbb{N}\dsep
        \forall u\in M\,\bigl((x,u)\in \RecCore{m_0}{f}\rightarrow
        \forall c\in M\,\bigl((x,c)\in \RecCore{m_0}{f}\rightarrow u=c\bigr)\bigr)\dsep
        (\Succ(x),d)\in \RecCore{m_0}{f}\dsep
        (\Succ(x),v)\in \RecCore{m_0}{f}\vdash d=v
      }{1,2,3,4,6,8}}

    \proofstepwidestar[1,2,3,5,6,7]{%
      (\Succ(x),c)\in \RecCore{m_0}{f}\rightarrow a=c
    }{%
      \rRI{8,9}}

    \proofstepwidestar[1,2,3,5,6]{%
      \forall c\in M\,\bigl((\Succ(x),c)\in \RecCore{m_0}{f}\rightarrow a=c\bigr)
    }{%
      \rUI{\rRI{7,10}}}

    \proofstepwidestar[1,2,3,5]{%
      \begin{aligned}[t]
        &(\Succ(x),a)\in \RecCore{m_0}{f}\rightarrow{}\\
        &\forall c\in M\,\bigl((\Succ(x),c)\in \RecCore{m_0}{f}\rightarrow a=c\bigr)
      \end{aligned}
    }{%
      \rRI{6,11}}

    \proofstepwidestar[1,2,3]{%
      \begin{aligned}[t]
        &\forall a\in M\,\bigl(
        (\Succ(x),a)\in \RecCore{m_0}{f}\rightarrow{}\\
        &\forall c\in M\,\bigl(
        (\Succ(x),c)\in \RecCore{m_0}{f}\rightarrow{}\\
        &a=c
        \bigr)\bigr)
      \end{aligned}
    }{%
      \rUI{\rRI{5,12}}}

    \proofstepwidestar[1,2,3,4]{%
      \begin{aligned}[t]
        &\forall a\in M\,\bigl(
        (\Succ(x),a)\in \RecCore{m_0}{f}\rightarrow{}\\
        &\forall c\in M\,\bigl(
        (\Succ(x),c)\in \RecCore{m_0}{f}\rightarrow{}\\
        &a=c
        \bigr)\bigr)
      \end{aligned}
    }{%
      \rEE{9,10,11,13}}
  \closeproofpartwideindR

  \proofpartwide{Schluss}
    \proofstepwidestar[1]{m_0\in M}{\rA}
    \proofstepwidestar[2]{f\colon M\to M}{\rA}
    \proofstepwidestar[3]{n\in\mathbb{N}}{\rA}
    \proofstepwidestar[4]{a\in M}{\rA}
    \proofstepwidestar[5]{b\in M}{\rA}
    \proofstepwidestar[6]{(n,a)\in \RecCore{m_0}{f}}{\rA}
    \proofstepwidestar[7]{(n,b)\in \RecCore{m_0}{f}}{\rA}

    \proofstepwidestar[1,2]{\forall u\in M\,\bigl((0,u)\in \RecCore{m_0}{f}\rightarrow \forall c\in M\,\bigl((0,c)\in \RecCore{m_0}{f}\rightarrow u=c\bigr)\bigr)}{%
      \FormulaRefAuto{m_0\in M\dsep f\colon M\to M \vdash \forall u\in M\,\bigl((0,u)\in \RecCore{m_0}{f}\rightarrow \forall c\in M\,\bigl((0,c)\in \RecCore{m_0}{f}\rightarrow u=c\bigr)\bigr)}{1,2}}
    \proofstepwidestar[1,2]{\forall x\in\mathbb N\,\bigl((\forall u\in M\,\bigl((x,u)\in \RecCore{m_0}{f}\rightarrow \forall c\in M\,\bigl((x,c)\in \RecCore{m_0}{f}\rightarrow u=c\bigr)\bigr))\rightarrow(\forall a\in M\,\bigl((\Succ(x),a)\in \RecCore{m_0}{f}\rightarrow \forall c\in M\,\bigl((\Succ(x),c)\in \RecCore{m_0}{f}\rightarrow a=c\bigr)\bigr))\bigr)}{%
      \FormulaRefAuto{  m_0\in M\dsep f\colon M\to M\dsep x\in\mathbb{N}\dsep
  \forall u\in M\,\bigl((x,u)\in \RecCore{m_0}{f}\rightarrow
  \forall c\in M\,\bigl((x,c)\in \RecCore{m_0}{f}\rightarrow u=c\bigr)\bigr)\vdash
  \forall a\in M\,\bigl((\Succ(x),a)\in \RecCore{m_0}{f}\rightarrow
  \forall c\in M\,\bigl((\Succ(x),c)\in \RecCore{m_0}{f}\rightarrow a=c\bigr)\bigr)}{1,2}}
    \proofstepwidestar[1,2,3]{%
      \begin{aligned}[t]
        &\forall a\in M\,\bigl((n,a)\in \RecCore{m_0}{f}\rightarrow{}\\
        &\qquad \forall c\in M\,\bigl((n,c)\in \RecCore{m_0}{f}\rightarrow a=c\bigr)\bigr)
      \end{aligned}
    }{\FormulaRefAuto{PeanoInductionRule}{8,9,3}}

    \proofstepwidestar[1,2,3,4]{%
      \begin{aligned}[t]
        &(n,a)\in \RecCore{m_0}{f}\rightarrow{}\\
        &\forall c\in M\,\bigl((n,c)\in \RecCore{m_0}{f}\rightarrow a=c\bigr)
      \end{aligned}
    }{%
      \rRE{\rUE{10},4}}

    \proofstepwidestar[1,2,3,4,6]{%
      \forall c\in M\,\bigl((n,c)\in \RecCore{m_0}{f}\rightarrow a=c\bigr)
    }{%
      \rRE{6,11}}

    \proofstepwidestar[1,2,3,4,5]{%
      (n,b)\in \RecCore{m_0}{f}\rightarrow a=b
    }{%
      \rRE{\rUE{12},5}}

    \proofstepwidestar[1,2,3,4,5,6,7]{a=b}{%
      \rRE{7,13}}
  \closeproofpartwide
\end{tabproofsplitwide}

\paragraph{Eindeutige Existenz}

\par\Needspace{14\baselineskip}
\FormulaThmDelta[Eindeutige Existenz eines Werts in jeder Stufe]{%
  m_0\in M\dsep f\colon M\to M\dsep n\in\mathbb{N}
  \vdash \exists! a\in M\,\bigl((n,a)\in \RecCore{m_0}{f}\bigr)
}{%
  \DeltaDecl{Mengen}{M\dsep n\dsep a\dsep b\dsep m_0}%
  \DeltaDecl{Funktionen}{f}%
}
\begin{tabproofsplit}
\proofpart{Existenz}
\proofstep{1}{m_0\in M}{\rA}
  \proofstep{2}{f\colon M\to M}{\rA}
  \proofstep{3}{n\in\mathbb{N}}{\rA}
  \proofstep{1,2,3}{\exists a\in M\,\bigl((n,a)\in \RecCore{m_0}{f}\bigr)}{%
    \FormulaRefAuto{m_0\in M\dsep f\colon M\to M\dsep n\in\mathbb{N}\vdash \exists a\in M\,\bigl((n,a)\in \RecCore{m_0}{f}\bigr)}{1,2,3}}

\closeproofpart

\proofpart{Eindeutigkeit}
  \setcounter{proofstepnr}{4}% Fortlaufende Schritte dieses Beweises.
\proofstep{5}{a\in M}{\rA}
  \proofstep{6}{b\in M}{\rA}
  \proofstep{7}{(n,a)\in \RecCore{m_0}{f}}{\rA}
  \proofstep{8}{(n,b)\in \RecCore{m_0}{f}}{\rA}
  \proofstep{1,2,3,5,6,7,8}{a=b}{%
    \FormulaRefAuto{m_0\in M\dsep f\colon M\to M\dsep n\in\mathbb{N}\dsep a\in M\dsep b\in M\dsep (n,a)\in \RecCore{m_0}{f}\dsep (n,b)\in \RecCore{m_0}{f}\vdash a=b}{1,2,3,5,6,7,8}}

\closeproofpart

\proofpart{Eindeutige Existenz}
  \setcounter{proofstepnr}{9}% Fortlaufende Schritte dieses Beweises.
\proofstep{1,2,3}{\exists! a\in M\,\bigl((n,a)\in \RecCore{m_0}{f}\bigr)}{%
    \UEI{4,5,6,7,8,9}}

\closeproofpart
\end{tabproofsplit}

\fi

\subsubsection{Der Rekursionskern als Funktion}

Nachdem jede Stufe genau einen Wert besitzt, kann der Rekursionskern selbst als
Graph einer Funktion \(\mathbb{N}\to M\) gelesen werden.

\paragraph{Anwendung auf den Rekursionskern}

\par\Needspace{14\baselineskip}
\FormulaThmDelta[Der Rekursionskern ist eine Abbildung]{%
  m_0\in M\dsep f\colon M\to M \vdash \RecCore{m_0}{f}\colon \mathbb{N}\to M
}{%
  \DeltaDecl{Mengen}{M\dsep x\dsep y\dsep z\dsep m_0}%
  \DeltaDecl{Funktionen}{f}%
}
\begin{tabproof}
  \proofstep{1}{m_0\in M}{\rA}
  \proofstep{2}{f\colon M\to M}{\rA}
  \proofstep{1,2}{\RecCore{m_0}{f}\subseteq \mathbb{N}\times M}{%
    \FormulaRefAuto{m_0\in M\dsep f\colon M\to M \vdash \RecCore{m_0}{f}\subseteq \mathbb{N}\times M}{1,2}}
  \proofstep{4}{n\in\mathbb{N}}{\rA}
  \proofstep{1,2,4}{%
    \exists! a\in M\,\bigl((n,a)\in \RecCore{m_0}{f}\bigr)
  }{\FormulaRefAuto{RecCoreStageUnique}{1,2,4}}
  \proofstep{1,2}{%
    \forall n\in\mathbb{N}\,
    \exists! a\in M\,\bigl((n,a)\in \RecCore{m_0}{f}\bigr)
  }{\rUI{\rRI{4,5}}}
  \proofstep{1,2}{\RecCore{m_0}{f}\colon \mathbb{N}\to M}{%
    \FormulaRefAuto{UniqueValuedGraphFunction}{3,6}}
\end{tabproof}

\paragraph{Rekursionsgleichungen}

\par\Needspace{14\baselineskip}
\FormulaThmDelta[Anfangswert des Rekursionskerns]{%
  m_0\in M\dsep f\colon M\to M \vdash \RecCore{m_0}{f}(0)=m_0
}{%
  \DeltaDecl{Mengen}{M\dsep m_0}%
  \DeltaDecl{Funktionen}{f}%
}
\begin{tabproof}
  \proofstep{1}{m_0\in M}{\rA}
  \proofstep{2}{f\colon M\to M}{\rA}

  \proofstep{1,2}{\RecCore{m_0}{f}\colon \mathbb{N}\to M}{%
    \FormulaRefAuto{m_0\in M\dsep f\colon M\to M \vdash \RecCore{m_0}{f}\colon \mathbb{N}\to M}{1,2}}

  \proofstep{1,2}{(0,m_0)\in \RecCore{m_0}{f}}{%
    \FormulaRefAuto{m_0\in M\dsep f\colon M\to M \vdash (0,m_0)\in \RecCore{m_0}{f}}{1,2}}

  \proofstep{1,2}{\RecCore{m_0}{f}(0)=m_0}{%
    \FormulaRefAuto{F\colon A\to B\dsep (x,y)\in F\vdash F(x)=y}{3,4}}
\end{tabproof}

\par\Needspace{14\baselineskip}
\FormulaThmDeltaR[Rekursionsgleichung des Rekursionskerns]{%
  \begin{aligned}
    &m_0\in M\dsep f\colon M\to M\dsep n\in\mathbb{N}\\
    &\vdash \RecCore{m_0}{f}(\Succ(n))=f(\RecCore{m_0}{f}(n))
  \end{aligned}
}{%
  m_0\in M\dsep f\colon M\to M\dsep n\in\mathbb{N}\vdash \RecCore{m_0}{f}(\Succ(n))=f(\RecCore{m_0}{f}(n))
}{%
  \DeltaDecl{Mengen}{M\dsep n\dsep m_0}%
  \DeltaDecl{Funktionen}{f}%
}
\begin{tabproof}
  \proofstep{1}{m_0\in M}{\rA}
  \proofstep{2}{f\colon M\to M}{\rA}
  \proofstep{3}{n\in\mathbb{N}}{\rA}

  \proofstep{1,2}{\RecCore{m_0}{f}\colon \mathbb{N}\to M}{%
    \FormulaRefAuto{m_0\in M\dsep f\colon M\to M \vdash \RecCore{m_0}{f}\colon \mathbb{N}\to M}{1,2}}

  \proofstep{1,2,3}{(n,\RecCore{m_0}{f}(n))\in \RecCore{m_0}{f}}{%
    \FormulaRefAuto{%
      F\colon A\to B\dsep x\in A\vdash (x,F(x))\in F
    }{4,3}}

  \proofstep{1,2,3}{\RecCore{m_0}{f}(n)\in M}{%
    \FormulaRefAuto{F\colon A\to B\dsep x\in A\vdash F(x)\in B}{4,3}}

  \proofstep{1,2,3}{(\Succ(n),f(\RecCore{m_0}{f}(n)))\in \RecCore{m_0}{f}}{%
    \FormulaRefAuto{m_0\in M\dsep f\colon M\to M\dsep n\in\mathbb{N}\dsep a\in M\dsep (n,a)\in \RecCore{m_0}{f}\vdash (\Succ(n),f(a))\in \RecCore{m_0}{f}}{1,2,3,6,5}}

  \proofstep{1,2,3}{\RecCore{m_0}{f}(\Succ(n))=f(\RecCore{m_0}{f}(n))}{%
    \FormulaRefAuto{%
      F\colon A\to B\dsep (x,y)\in F\vdash F(x)=y
    }{4,7}}
\end{tabproof}

\paragraph{Eindeutigkeit rekursiv definierter Abbildungen}

Die folgende Induktion trennt die spätere Eindeutigkeitsaussage vom
Existenzbeweis: Jede Abbildung mit demselben Anfangswert und derselben
Rekursionsgleichung stimmt punktweise mit jeder anderen solchen Abbildung
überein.

\par\Needspace{14\baselineskip}
\FormulaThmDeltaR[Eindeutigkeit rekursiv definierter Abbildungen]{%
  \begin{aligned}
    &G\colon \mathbb{N}\to M\dsep H\colon \mathbb{N}\to M\dsep
    G(0)=m_0\dsep H(0)=m_0\dsep{}\\
    &\forall n\in\mathbb{N}\,G(\Succ(n))=f(G(n))\dsep
    \forall n\in\mathbb{N}\,H(\Succ(n))=f(H(n))
    \vdash G=H
  \end{aligned}
}{%
  G\colon \mathbb{N}\to M\dsep H\colon \mathbb{N}\to M\dsep
  G(0)=m_0\dsep H(0)=m_0\dsep
  \forall n\in\mathbb{N}\,G(\Succ(n))=f(G(n))\dsep
  \forall n\in\mathbb{N}\,H(\Succ(n))=f(H(n))
  \vdash G=H
}{%
  \DeltaDecl{Mengen}{M\dsep n\dsep x\dsep m_0}%
  \DeltaDecl{Funktionen}{f\dsep G\dsep H}%
}
\begin{tabproofsplitwide}
  \proofpartwideindR[Induktionsanfang]{%
    \begin{aligned}[t]
      &G\colon \mathbb{N}\to M\dsep H\colon \mathbb{N}\to M\dsep{}\\
      &G(0)=m_0\dsep H(0)=m_0\dsep{}\\
      &\forall n\in\mathbb{N}\,G(\Succ(n))=f(G(n))\dsep{}\\
      &\forall n\in\mathbb{N}\,H(\Succ(n))=f(H(n))
      \vdash G(0)=H(0)
    \end{aligned}
  }{%
    G\colon \mathbb{N}\to M\dsep H\colon \mathbb{N}\to M\dsep
    G(0)=m_0\dsep H(0)=m_0\dsep
    \forall n\in\mathbb{N}\,G(\Succ(n))=f(G(n))\dsep
    \forall n\in\mathbb{N}\,H(\Succ(n))=f(H(n))
    \vdash G(0)=H(0)
  }
    \proofstepwidestar[1]{G\colon \mathbb{N}\to M}{\rA}
    \proofstepwidestar[2]{H\colon \mathbb{N}\to M}{\rA}
    \proofstepwidestar[3]{G(0)=m_0}{\rA}
    \proofstepwidestar[4]{H(0)=m_0}{\rA}
    \proofstepwidestar[5]{\forall n\in\mathbb{N}\,G(\Succ(n))=f(G(n))}{\rA}
    \proofstepwidestar[6]{\forall n\in\mathbb{N}\,H(\Succ(n))=f(H(n))}{\rA}

    \proofstepwidestar[3,4]{G(0)=H(0)}{%
      \FormulaRefAuto{a = b,\, c = b \vdash a = c}{3,4}}
  \closeproofpartwideindR

  \proofpartwideindR[Induktionsschritt]{%
    \begin{aligned}[t]
      &G\colon \mathbb{N}\to M\dsep H\colon \mathbb{N}\to M\dsep{}\\
      &G(0)=m_0\dsep H(0)=m_0\dsep{}\\
      &\forall n\in\mathbb{N}\,G(\Succ(n))=f(G(n))\dsep{}\\
      &\forall n\in\mathbb{N}\,H(\Succ(n))=f(H(n))\dsep{}\\
      &x\in\mathbb{N}\dsep G(x)=H(x)\vdash G(\Succ(x))=H(\Succ(x))
    \end{aligned}
  }{%
    G\colon \mathbb{N}\to M\dsep H\colon \mathbb{N}\to M\dsep
    G(0)=m_0\dsep H(0)=m_0\dsep
    \forall n\in\mathbb{N}\,G(\Succ(n))=f(G(n))\dsep
    \forall n\in\mathbb{N}\,H(\Succ(n))=f(H(n))\dsep
    x\in\mathbb{N}\dsep G(x)=H(x)\vdash G(\Succ(x))=H(\Succ(x))
  }
    \proofstepwidestar[1]{G\colon \mathbb{N}\to M}{\rA}
    \proofstepwidestar[2]{H\colon \mathbb{N}\to M}{\rA}
    \proofstepwidestar[3]{G(0)=m_0}{\rA}
    \proofstepwidestar[4]{H(0)=m_0}{\rA}
    \proofstepwidestar[5]{\forall n\in\mathbb{N}\,G(\Succ(n))=f(G(n))}{\rA}
    \proofstepwidestar[6]{\forall n\in\mathbb{N}\,H(\Succ(n))=f(H(n))}{\rA}
    \proofstepwidestar[7]{x\in\mathbb{N}}{\rA}
    \proofstepwidestar[8]{G(x)=H(x)}{\rA}

    \proofstepwidestar[1,7]{G(x)\in M}{%
      \FormulaRefAuto{F\colon A\to B\dsep x\in A\vdash F(x)\in B}{1,7}}

    \proofstepwidestar[2,7]{H(x)\in M}{%
      \FormulaRefAuto{F\colon A\to B\dsep x\in A\vdash F(x)\in B}{2,7}}

    \proofstepwidestar[5]{x\in\mathbb{N}\rightarrow G(\Succ(x))=f(G(x))}{\rUE{5}}
    \proofstepwidestar[5,7]{G(\Succ(x))=f(G(x))}{\rRE{11,7}}

    \proofstepwidestar[6]{x\in\mathbb{N}\rightarrow H(\Succ(x))=f(H(x))}{\rUE{6}}
    \proofstepwidestar[6,7]{H(\Succ(x))=f(H(x))}{\rRE{13,7}}

    \proofstepwidestar[1,2,7,8]{f(G(x))=f(H(x))}{%
      \rIE{8,\rII}}

    \proofstepwidestar[1,2,6,7,8]{H(\Succ(x))=f(G(x))}{%
      \FormulaRefAuto{a = b,\, c = b \vdash a = c}{14,15}}

    \proofstepwidestar[1,2,6,7,8]{f(G(x))=H(\Succ(x))}{%
      \FormulaRefAuto{a=b\vdash b=a}{16}}

    \proofstepwidestar[1,2,5,6,7,8]{G(\Succ(x))=H(\Succ(x))}{%
      \FormulaRefAuto{a = b,\, c = b \vdash a = c}{12,17}}
  \closeproofpartwideindR

  \proofpartwide{Schluss}
    \proofstepwidestar[1]{G\colon \mathbb{N}\to M}{\rA}
    \proofstepwidestar[2]{H\colon \mathbb{N}\to M}{\rA}
    \proofstepwidestar[3]{G(0)=m_0}{\rA}
    \proofstepwidestar[4]{H(0)=m_0}{\rA}
    \proofstepwidestar[5]{\forall n\in\mathbb{N}\,G(\Succ(n))=f(G(n))}{\rA}
    \proofstepwidestar[6]{\forall n\in\mathbb{N}\,H(\Succ(n))=f(H(n))}{\rA}

    \proofstepwidestar[7]{n\in\mathbb{N}}{\rA}

    \proofstepwidestar[1,2,3,4,5,6]{G(0)=H(0)}{%
      \FormulaRefAuto{%
          G\colon \mathbb{N}\to M\dsep H\colon \mathbb{N}\to M\dsep
          G(0)=m_0\dsep H(0)=m_0\dsep
          \forall n\in\mathbb{N}\,G(\Succ(n))=f(G(n))\dsep
          \forall n\in\mathbb{N}\,H(\Succ(n))=f(H(n))
          \vdash G(0)=H(0)
        }{1,2,3,4,5,6}}
    \proofstepwidestar[1,2,3,4,5,6]{\forall x\in\mathbb N\,\bigl((G(x)=H(x))\rightarrow(G(\Succ(x))=H(\Succ(x)))\bigr)}{%
      \FormulaRefAuto{%
          G\colon \mathbb{N}\to M\dsep H\colon \mathbb{N}\to M\dsep
          G(0)=m_0\dsep H(0)=m_0\dsep
          \forall n\in\mathbb{N}\,G(\Succ(n))=f(G(n))\dsep
          \forall n\in\mathbb{N}\,H(\Succ(n))=f(H(n))\dsep
          x\in\mathbb{N}\dsep G(x)=H(x)\vdash G(\Succ(x))=H(\Succ(x))
        }{1,2,3,4,5,6}}
    \proofstepwidestar[1,2,3,4,5,6,7]{G(n)=H(n)}{\FormulaRefAuto{PeanoInductionRule}{8,9,7}}

    \proofstepwidestar[1,2,3,4,5,6]{\forall n\in\mathbb{N}\,\bigl(G(n)=H(n)\bigr)}{%
      \rUI{\rRI{7,10}}}

    \proofstepwidestar[1,2,3,4,5,6]{G=H}{%
      \FormulaRefAuto{FunctionExtensionality}[thm]{1,2,11}}
  \closeproofpartwide
\end{tabproofsplitwide}

\subsubsection{Der Dedekindsche Rekursionssatz}

Der Hauptsatz ist nun nur noch die Zusammenführung der beiden vorherigen
Blöcke. Die Existenz kommt vom Rekursionskern als Funktion, die Eindeutigkeit
von der punktweisen Induktion über rekursiv definierte Abbildungen.

\par\Needspace{14\baselineskip}
\FormulaThmDeltaR[Existenz einer rekursiv definierten Abbildung]{%
  \begin{aligned}
    &m_0\in M\dsep f\colon M\to M\\
    \vdash{}&\exists G\colon \mathbb{N}\to M\,\Bigl(
      G(0)=m_0 \land{}\\
      &\qquad\qquad\qquad\forall n\in\mathbb{N}\,
      G(\Succ(n))=f(G(n))
    \Bigr)
  \end{aligned}
}{%
  m_0\in M\dsep f\colon M\to M\vdash \exists G\colon \mathbb{N}\to M\,\Bigl(G(0)=m_0 \land \forall n\in\mathbb{N}\,G(\Succ(n))=f(G(n))\Bigr)
}{%
  \DeltaDecl{Mengen}{M\dsep n\dsep m_0}%
  \DeltaDecl{Funktionen}{f\dsep G}%
}
\begin{tabproofsplit}
\proofpart{Funktionstyp}
\proofstep{1}{m_0\in M}{\rA}
  \proofstep{2}{f\colon M\to M}{\rA}

  \proofstep{1,2}{\RecCore{m_0}{f}\colon \mathbb{N}\to M}{%
    \FormulaRefAuto{m_0\in M\dsep f\colon M\to M \vdash \RecCore{m_0}{f}\colon \mathbb{N}\to M}{1,2}}


\closeproofpart

\proofpart{Anfangswert}
  \setcounter{proofstepnr}{3}% Fortlaufende Schritte dieses Beweises.
\proofstep{1,2}{\RecCore{m_0}{f}(0)=m_0}{%
    \FormulaRefAuto{m_0\in M\dsep f\colon M\to M \vdash \RecCore{m_0}{f}(0)=m_0}{1,2}}


\closeproofpart

\proofpart{Rekursionsgleichung}
  \setcounter{proofstepnr}{4}% Fortlaufende Schritte dieses Beweises.
\proofstep{5}{n\in\mathbb{N}}{\rA}

  \proofstep{1,2,5}{%
    \begin{aligned}[t]
      &\RecCore{m_0}{f}(\Succ(n))={}\\
      &f(\RecCore{m_0}{f}(n))
    \end{aligned}
  }{%
    \FormulaRefAuto{m_0\in M\dsep f\colon M\to M\dsep n\in\mathbb{N}\vdash \RecCore{m_0}{f}(\Succ(n))=f(\RecCore{m_0}{f}(n))}{1,2,5}}

  \proofstep{1,2}{%
    \begin{aligned}[t]
      &\forall n\in\mathbb{N}\,\bigl({}\\
      &\qquad \RecCore{m_0}{f}(\Succ(n))={}\\
      &\qquad f(\RecCore{m_0}{f}(n))\bigr)
    \end{aligned}
  }{%
    \rUI{\rRI{5,6}}}


\closeproofpart

\proofpart{Zusammenführung und Existenz}
  \setcounter{proofstepnr}{7}% Fortlaufende Schritte dieses Beweises.
\proofstep{1,2}{%
    \begin{aligned}[t]
      &\RecCore{m_0}{f}(0)=m_0 \land{}\\
      &\forall n\in\mathbb{N}\,\bigl({}\\
      &\qquad \RecCore{m_0}{f}(\Succ(n))={}\\
      &\qquad f(\RecCore{m_0}{f}(n))\bigr)
    \end{aligned}
  }{\rAI{4,7}}

  \proofstep{1,2}{%
    \begin{aligned}[t]
      &\exists G\colon \mathbb{N}\to M\,\Bigl({}\\
      &\qquad G(0)=m_0 \land{}\\
      &\qquad \forall n\in\mathbb{N}\,\bigl({}\\
      &\qquad\qquad G(\Succ(n))=f(G(n))\bigr)\Bigr)
    \end{aligned}
  }{%
    \rEI{\rAI{3,8}}}

\closeproofpart
\end{tabproofsplit}



\par\Needspace{12\baselineskip}
\hypertarget{dedekind.proof}{}
\noindent\textbf{Beweis des Dedekindschen Rekursionssatzes.}\par\nobreak
\noindent Hauptsatz: \FormulaRefAuto{DedekindRecursionTheorem}[thm] in Band~10.
\par\medskip

\begin{tabproofsplit}
\proofpart{Existenz}
\proofstep{1}{m_0\in M}{\rA}
  \proofstep{2}{f\colon M\to M}{\rA}

  \proofstep{1,2}{%
    \begin{aligned}[t]
      &\exists G\colon \mathbb{N}\to M\,\Bigl({}\\
      &\qquad G(0)=m_0 \land{}\\
      &\qquad \forall n\in\mathbb{N}\,\bigl({}\\
      &\qquad\qquad G(\Succ(n))=f(G(n))\bigr)\Bigr)
    \end{aligned}
  }{%
    \FormulaRefAuto{m_0\in M\dsep f\colon M\to M\vdash \exists G\colon \mathbb{N}\to M\,\Bigl(G(0)=m_0 \land \forall n\in\mathbb{N}\,G(\Succ(n))=f(G(n))\Bigr)}{1,2}}


\closeproofpart

\proofpart{Eindeutigkeit}
  \setcounter{proofstepnr}{3}% Fortlaufende Schritte dieses Beweises.
\proofstep{4}{G\colon \mathbb{N}\to M}{\rA}

  \proofstep{5}{%
    \begin{aligned}[t]
      &G(0)=m_0 \land{}\\
      &\forall n\in\mathbb{N}\,\bigl({}\\
      &\qquad G(\Succ(n))=f(G(n))\bigr)
    \end{aligned}
  }{\rA}

  \proofstep{6}{H\colon \mathbb{N}\to M}{\rA}

  \proofstep{7}{%
    \begin{aligned}[t]
      &H(0)=m_0 \land{}\\
      &\forall n\in\mathbb{N}\,\bigl({}\\
      &\qquad H(\Succ(n))=f(H(n))\bigr)
    \end{aligned}
  }{\rA}

  \proofstep{5}{G(0)=m_0}{\rAEa{5}}

  \proofstep{5}{%
    \begin{aligned}[t]
      &\forall n\in\mathbb{N}\,\bigl({}\\
      &\qquad G(\Succ(n))=f(G(n))\bigr)
    \end{aligned}
  }{\rAEb{5}}

  \proofstep{7}{H(0)=m_0}{\rAEa{7}}

  \proofstep{7}{%
    \begin{aligned}[t]
      &\forall n\in\mathbb{N}\,\bigl({}\\
      &\qquad H(\Succ(n))=f(H(n))\bigr)
    \end{aligned}
  }{\rAEb{7}}

  \proofstep{4,5,6,7}{G=H}{%
    \FormulaRefAuto{G\colon \mathbb{N}\to M\dsep H\colon \mathbb{N}\to M\dsep G(0)=m_0\dsep H(0)=m_0\dsep \forall n\in\mathbb{N}\,G(\Succ(n))=f(G(n))\dsep \forall n\in\mathbb{N}\,H(\Succ(n))=f(H(n))\vdash G=H}{4,6,8,10,9,11}}


\closeproofpart

\proofpart{Eindeutige Existenz}
  \setcounter{proofstepnr}{12}% Fortlaufende Schritte dieses Beweises.
\proofstep{1,2}{%
    \begin{aligned}[t]
      &\exists! G\colon \mathbb{N}\to M\,\Bigl({}\\
      &\qquad G(0)=m_0 \land{}\\
      &\qquad \forall n\in\mathbb{N}\,\bigl({}\\
      &\qquad\qquad G(\Succ(n))=f(G(n))\bigr)\Bigr)
    \end{aligned}
  }{%
    \UEI{3,4,5,6,7,12}}

\closeproofpart
\end{tabproofsplit}


\paragraph{Anschluss an die Rekursionsabbildung}
Der Rekursionskern ist eine Funktion mit den beiden verlangten
Rekursionsgleichungen. Nach der bewiesenen Eindeutigkeit stimmt er daher
mit der in \FormulaRefAuto{RecFunDef}[def] benannten Funktion überein:
\[
  m_0\in M,\quad f\colon M\to M
  \quad\Longrightarrow\quad
  \RecCore{m_0}{f}=\RecFun{m_0}{f}.
\]
Die Definition im Hauptband setzt dadurch keine Kenntnis der
Schnittkonstruktion voraus.
